\documentclass[usenatbib,twoside,twocolumn]{mnras}

% Paper packages
\usepackage{graphicx}
\usepackage{amsmath}
\usepackage{amssymb}
\usepackage{booktabs}

% Denser float packing (reduce float-only pages; no content change)
\renewcommand{\topfraction}{0.9}
\renewcommand{\bottomfraction}{0.85}
\renewcommand{\textfraction}{0.07}
\renewcommand{\floatpagefraction}{0.75}
\renewcommand{\dbltopfraction}{0.9}
\renewcommand{\dblfloatpagefraction}{0.75}
\setcounter{topnumber}{3}
\setcounter{bottomnumber}{2}
\setcounter{totalnumber}{5}
\setcounter{dbltopnumber}{3}

% Define astronomical units
\newcommand{\msun}{\ensuremath{M_{\odot}}}
\newcommand{\pc}{\ensuremath{\rm pc}}
\newcommand{\kms}{\ensuremath{\rm km\,s^{-1}}}

\title[Filament fragmentation and core spacing -- I]{Filament fragmentation and core spacing -- I.\\ Two spacing statistics and a clustered core distribution in \textit{Herschel} Gould Belt clouds}

\author[G.~J.~White]{G.~J.~White$^{1,2}$\\
$^{1}$School of Physical Sciences, The Open University, Walton Hall, Milton Keynes, MK7 6AA, UK\\
$^{2}$Rutherford Appleton Laboratory, Harwell Science and Innovation Campus, Didcot, OX11 0QX, UK}

\date{\today}
\pubyear{2026}

\begin{document}
\label{firstpage}
\maketitle

\begin{abstract}

Dense cores on molecular-cloud filaments are widely reported at separations shorter than the fastest-growing mode of an equilibrium isothermal cylinder, $\lambda_{\rm m}\approx5$--$6\,W$. We ask what that comparison measures. We re-measure core separations along filament crests in four \textit{Herschel} Gould Belt clouds using \textit{Gaia} DR3 distances, and separate two statistics that the literature conflates: the adjacent-core spacing within populated stretches, $S_{\rm local}$, and the traced skeleton length per associated core, $S_{\rm skel}$.

They disagree. Under the adopted persistent-crest network, $0.06$\,pc association radius and fixed width, $S_{\rm local}=1.7$--$3.0$ and $S_{\rm skel}=3.4$--$5.8$; the second lies near the equilibrium-cylinder reference in two clouds, but not under other defensible choices, whose envelope spans a factor of eight. Extracting the network as a mask medial axis rather than as persistent crests moves $S_{\rm local}$ by a factor of two. Cores occupy only $6$--$27$ per cent of the crest, in groups separated by long empty stretches, and their gap distributions reject a periodic comb, a deliberately imperfect comb and a homogeneous Poisson process, showing no feature at the equilibrium-cylinder scale.

We conclude that no single characteristic spacing describes these cores, that a fragmentation model must reproduce a spatial point process rather than a wavelength, and that the limiting uncertainty is the structural definition of a filament, not measurement precision. We also show that for points distributed uniformly along, or regularly across, a segment the all-pairs median converges on $0.293\,L$; more generally it behaves as an extent statistic rather than a local-spacing estimator. Paper~II asks whether classical fragmentation theory need apply at all.

\end{abstract}

\begin{keywords}
stars: formation -- ISM: structure -- ISM: filaments -- ISM: clouds
\end{keywords}

\section{Introduction}
\label{sec:intro}

The \textit{Herschel} Space Observatory established the ubiquity of filamentary structure in nearby star-forming clouds \citep{Andre2010} and a characteristic width of $\sim$0.1\,pc across diverse environments \citep{Arzoumanian2011}, though the universality and physical interpretation of that scale remain contested \citep{Panopoulou2017,Panopoulou2022}. For an isothermal cylinder the critical mass per unit length is $M_{\rm line,crit}\approx16\,M_\odot\,{\rm pc^{-1}}$ \citep{Ostriker1964}, and a marginally-stable cylinder fragments at $\lambda_{\rm m}\approx5$--$6\times$FWHM \citep{Inutsuka1992,Fischera2012}. Observed core separations are widely reported to fall below that value, and the deficit has been read as evidence for additional physics: hierarchical fragmentation into velocity-coherent fibres \citep{Hacar2013,Hacar2018,Yang2024}, magnetic tension \citep{Nakamura1993}, or fragmentation from a dynamically evolving rather than an equilibrium state \citep{Andre2014,Clarke2016,Clarke2017,Clarke2020}.

This paper asks a deliberately simple question. \emph{How far apart are the cores on HGBS filaments, and is that consistent with the equilibrium-cylinder calculation?} We find that the question does not have a single answer, and that the reason it does not is itself the result.

\smallskip\noindent\textbf{Three numbers.} Everything below rests on three quantities, and we introduce no others in the main text.
\begin{itemize}\setlength{\itemsep}{1pt}
\item $s_{\rm ACS}$, the median separation between \emph{neighbouring} cores measured along the crest. This is a conditional quantity: it is measured only where cores occur, and says nothing about filament that has formed none.
\item $L_{\rm skel}/N$, the traced skeleton length per associated core. This is unconditional. It is the closest \emph{available} analogue of a model that makes one fragment per wavelength, but it is not that model's observable: the theory refers to a coherent, supercritical, self-gravitating cylinder, whereas the traced skeleton also contains subcritical stretches, low-contrast extensions and material that need not be connected in velocity. We therefore treat it as an operational line-density statistic and not as an empirical fragmentation wavelength.
\item $f_{\rm occ}$, the fraction of crest length lying within a chosen distance of a source. It is a scale-dependent coverage statistic and not an intrinsic filling factor: we quote it at $\pm0.06$\,pc, one half-width, and give the full curve $f_{\rm occ}(\ell)$ where the scale matters.
\end{itemize}
We write $S_{\rm local}=s_{\rm ACS}/W$ and $S_{\rm skel}=L_{\rm skel}/(NW)$ for these in units of the filament width $W$, so that they can be compared with $\lambda_{\rm m}/W\approx5$--$6$. Table~\ref{tab:notation} fixes the notation, and Fig.~\ref{fig:concept} shows why the two statistics differ and what their difference measures.

\begin{table}
\caption{Notation. The distinction between a \emph{separation} (a measured distance between two objects), a \emph{line density} (objects per unit length) and a \emph{wavelength} (a property of an instability mode) is maintained throughout; the three are not interchangeable.}
\label{tab:notation}
\footnotesize
\setlength{\tabcolsep}{4pt}
\begin{tabular}{@{}lp{6.0cm}@{}}
\toprule
Symbol & Meaning \\
\midrule
$s_{\rm ACS}$ & median separation between along-crest neighbouring associated sources (a separation) \\
$L_{\rm skel}$ & total traced skeleton length \\
$N$ & number of catalogued sources associated with the skeleton \\
$W$ & filament width; fixed $0.10$\,pc unless the region-Plummer FWHM is specified \\
$S_{\rm local}$ & $s_{\rm ACS}/W$: conditional, measured only within core-bearing stretches \\
$S_{\rm skel}$ & $L_{\rm skel}/(NW)$: unconditional line density in width units \\
$f_{\rm occ}(\ell)$ & fraction of skeleton length within $\ell$ of a source; quoted at $\ell=0.06$\,pc \\
$S_{\rm elig}$ & exploratory: $S_{\rm skel}$ restricted to nominally supercritical crest \\
$\lambda_{\rm m}$ & fastest-growing wavelength of the equilibrium cylinder, $\approx5$--$6\,W$ (a wavelength) \\
$\lambda_{\rm cr}$ & classical critical wavelength, $\approx2.65\,W$ \\
\bottomrule
\end{tabular}
\end{table}

\smallskip\noindent\textbf{What we find.} $S_{\rm local}=1.7$--$3.0$ and $S_{\rm skel}=3.4$--$5.8$ under the adopted conventions. The gap between them is a measurement rather than a discrepancy to be resolved: cores occupy only $6$--$27$ per cent of the crest, so a statistic conditioned on their presence and one averaged over the whole filament are asking different questions. On the unconditional statistic, two of the four clouds sit within the equilibrium-cylinder band and the other two fall below it by less than a factor of two. On the conditional statistic all four are several times below it. Neither number is a measurement of a fragmentation wavelength, and a model that predicts one wavelength cannot be tested against either without first specifying how instability modes map onto an observed point process.

\smallskip\noindent\textbf{Two things that limit the answer.} The first is the estimator. The all-pairs median separation, which appears in some of the spacing literature, is not an estimator of local spacing. For points distributed uniformly along, or spaced regularly across, a segment of length $L$ it converges on $0.293L$; more generally it follows the extent of the occupied region rather than the separation between neighbours, with a coefficient that depends on the distribution (Section~\ref{sec:statistics}). We use adjacent along-crest separations and retain the all-pairs statistic only for comparison with earlier work. The published HGBS values we could check are local or along-fibre statistics and are consistent with ours to within about a third in each cloud, so this is a caution for future work and not a correction to past work.

The second is more fundamental. The equilibrium-cylinder calculation is an eigenmode of a single coherent, approximately cylindrical, self-gravitating object. A skeleton traced through a projected column-density map need not be that object: it can join material unrelated in velocity, or divide material that is dynamically single. The consequence is quantitative. Extracting the network as persistent crests rather than as the medial axis of a mask changes $S_{\rm local}$ by a factor of two and $S_{\rm skel}$ by more (Section~\ref{sec:ontology}). That single choice dominates every other systematic in this paper combined, which is why the comparisons below are diagnostics and not tests.

\smallskip\noindent\textbf{Scope.} Four clouds are four independent units. All $\pm$ values quoted here are descriptive $N=4$ ranges and not confidence intervals, and nothing below is offered as a universal fragmentation law. The MHD simulations that ask whether an evolving filament need select the equilibrium wavelength are presented separately in Paper~II; the present paper is observational throughout.

\begin{figure*}
\includegraphics[width=0.78\textwidth]{figures/fig_concept.pdf}
\caption{Why one spacing statistic is not enough. Red circles are cores; the horizontal bar is the filament, blue where it is unstable, grey where it is not, and green where it is core-bearing on a $\pm$half-width definition. \emph{(a)} Periodic fragmentation over the whole unstable length: the adjacent-core spacing and the length per core coincide. \emph{(b)} The same ten cores gathered into three groups: $s_{\rm ACS}$ falls by a factor of three while $L/N$ is unchanged, and their ratio measures how unevenly the cores occupy the crest. \emph{(c)} The same cores on a longer skeleton that includes material which was never unstable: $s_{\rm ACS}$ and $L_{\rm skel}/N$ are both unaltered, but the physically relevant $L_{\rm unstable}/N$ is smaller than either, so the comparison with an equilibrium-cylinder wavelength depends on which length is used. Panels (b) and (c) between them describe the observations reported here.}
\label{fig:concept}
\end{figure*}

\section{Data and method}
\label{sec:data}

Our sample comprises 8 high-confidence HGBS regions (four \emph{primary}, Table~\ref{tab:sample}; four \emph{limited}, Appendix~\ref{sec:limited_sample}). The original catalogues \citep{Konyves2015, Konyves2020, Ladjelate2020, Pezzuto2021} adopted distances of $130$--$400$\,pc from the best estimates available in 2010--2015 (Table~\ref{tab:literature} lists the pre-Gaia values; note that \citealt{Konyves2020} already adopted $\sim400$\,pc for Orion~B); Gaia DR3 parallaxes \citep{Zhang2023} have since enabled substantial revisions, largest for the more distant regions, whose pre-Gaia distances were the most uncertain.

We adopt the Gaia DR3 YSO-based distances from \citet{Zhang2023} (Table~4), with typical uncertainties of 3\%--5\%. The revisions are strongly region-dependent (largest for Serpens $260\to458$\,pc, $+76\%$, and Aquila $260\to436$\,pc, $+68\%$; Orion~B, Taurus and Ophiuchus essentially unchanged); since spacings scale linearly with distance, the core-count-weighted mean rises only modestly ($\sim10\%$ for the primary regions), dominated by Aquila rather than being a uniform $+32\%$ shift. Because Gaia DR3 thus compresses or expands each cloud non-uniformly, the revision is not a single global rescaling, which is a further reason the clouds must be treated as $N=4$ independent units and not by pooling cores that carry very different distance corrections. The observed spacing is a line-density statistic, not a resonance, and coincides with the theoretical wavelength only to order of magnitude.

The analysis includes all HGBS-identified cores (starless, prestellar and protostellar), extracted and classified within the consistent HGBS framework \citep{Konyves2015, Konyves2020}. Including unbound starless and protostellar sources may introduce some bias, but fragmentation produces a continuum of boundness states, and restricting the sample to prestellar cores would shrink several regions (notably TMC1, CRA and Serpens) below statistical usefulness. Core positions and properties are taken directly from the published catalogues.

Protostellar migration displacements of $0.01$--$0.05$\,pc \citep{Kirk2016, Mattern2018} are small against the $\sim0.2$--$0.3$\,pc spacing; we estimate the resulting ACS bias at $\sim5$--$10\%$ and carry it as a systematic (Section~\ref{sec:systematics}).




\begin{table}
\caption{Primary HGBS sample with Gaia DR3 distances. Per-region spacings are pairwise-median values rounded to three significant figures; the core-count-weighted mean of the four primary regions is $0.290$\,pc. The four \emph{limited} regions (Ophiuchus, Serpens, TMC1, CRA) and the full 8-region weighted mean ($0.279$\,pc) are listed in Appendix~\ref{sec:limited_sample}.}
\label{tab:sample}
\footnotesize
\setlength{\tabcolsep}{3pt}
\begin{tabular}{lccccc}
\toprule
Region & $D$ & $N_{\rm core}$ & Pairwise med. & S.E. & Status \\
       & (pc) &  & (pc) & (pc) &  \\
\midrule
Orion B  & 386 & 1,870 & 0.313 & 0.047 & Primary \\
Aquila   & 436 & 749   & 0.346 & 0.047 & Primary \\
Perseus  & 300 & 816   & 0.248 & 0.040 & Primary \\
Taurus   & 135 & 536   & 0.198 & 0.040 & Primary \\
\midrule
\textbf{Weighted mean}$^\dagger$ &  & \textbf{3,971} & \textbf{0.290} & &  \\
\bottomrule
\end{tabular}
\par\vspace{2pt}
\footnotesize $^\dagger$Core-count-weighted mean of the four primary regions. $N_{\rm core}$ is the number of deposited catalogue entries with usable coordinates, which for Orion~B ($1870$) differs slightly from the $1844$ quoted by \citet{Konyves2020}; we use the deposited count throughout for internal consistency. A formal core-count-weighted standard error (treating cores as independent) is \emph{not} the relevant uncertainty, since the four clouds are the independent units (Section~\ref{sec:systematics}); the cloud-to-cloud scatter is carried instead.
\end{table}

\smallskip\noindent\textbf{A note on terminology.} The HGBS catalogues contain unbound starless sources, candidate prestellar cores and protostellar sources. Where the evolutionary state is immaterial to the argument, which is most of this paper, we write \emph{catalogued source} or simply \emph{source}, and reserve \emph{core} for the classified populations used in the population splits of Section~\ref{sec:coverage}. All the headline statistics use the all-source catalogue; Table~\ref{tab:popsplit} gives them for each classified subsample.

\subsection{Sample selection}
\label{sec:distance_uncertainties}

Three criteria, fixed before any spacing was measured, divide the eight regions into a \emph{primary} and a \emph{limited} sample (Table~\ref{tab:selection}): a catalogue of more than $500$ cores; a well-established distance, taken as a \textit{Gaia} DR3 uncertainty below $10$ per cent; and a skeleton simple enough for unambiguous core-to-filament association, taken as a largest connected component carrying at least a quarter of the total length. Four regions satisfy all three. Ophiuchus fails only the topology criterion, its skeleton being a dense, heavily branched network; TMC1 and CRA fail only the core-count criterion; and Serpens fails the distance criterion, lying within the Aquila/Serpens Rift, whose line-of-sight depth makes a single assigned distance ambiguous. The limited regions are excluded from the adjacent-core measurement and retained only for the legacy pairwise comparison.

 The largest distance revisions (Serpens $+76\%$, Aquila $+68\%$) are independently corroborated by VLBI maser parallaxes (validating Gaia DR3 to $<5\%$; \citealt{Kounkel2018,Xu2021,OrtizLeon2018,OrtizLeon2017}), extinction estimates \citep{Yan2022,Green2024} and co-spatiality with Aquila \citep{Zucker2020}; Serpens is \emph{limited} and excluded in any case. The short-spacing result is independent of these regions: leave-one-out shifts the weighted mean by $<6\%$, and dropping Aquila leaves the region-Plummer median at $1.40$ and the fixed-width characteristic at $2.1$, both well below the reference band. The split is not tuned to the distance revision: its criteria act on the distance-independent core count and skeleton morphology, the closest any region lies to the boundary is Taurus at $536$ cores, and because the \textit{Gaia} revision rescales spacings but not core counts no plausible distance change moves a region across the classification.

\smallskip\noindent\textbf{Are the thresholds themselves neutral?} They are operational and the fact that they were fixed in advance does not make them scientifically innocuous. The topology criterion in particular could preferentially admit clouds in which a one-dimensional analysis is easy and exclude the branched networks where a spacing description might fail worst. We therefore relaxed them and re-ran the full pipeline on the excluded clouds. Dropping the topology criterion admits Ophiuchus, whose $513$ catalogued sources give $411$ associations on only $32$\,pc of crest and yield $S_{\rm local}=0.6$ and $S_{\rm skel}=0.8$; lowering the core count to $150$ admits CRA, which gives $S_{\rm local}=2.7$ and $S_{\rm skel}=4.4$. TMC1's deposited skeleton carries no usable astrometric header and cannot be run through the same pipeline; it is in any case a sub-region of Taurus. Admitting the harder clouds therefore widens the envelope \emph{downwards}: Ophiuchus, the most branched network in the survey, is the furthest of all from the equilibrium-cylinder reference on both statistics, which is the opposite of a selection favouring agreement. It also makes the segmentation problem sharper rather than softer, because a heavily branched medial-axis-like crest counts short spurs as filament length. We report these values with quality flags in Appendix~\ref{sec:limited_sample} rather than dropping the clouds silently, but we do not fold them into the headline numbers, because the association step on such a network is not equivalent to the one performed on the primary four.


\begin{table}
\caption{Sample selection. All three criteria were fixed before any spacing statistic was computed. $N_{\rm core}$ is the published catalogue size; ``topology'' requires the largest connected skeleton component to carry at least a quarter of the total length. No core count is entered for Serpens, which is excluded on distance and for which we do not use a published catalogue.}
\label{tab:selection}
\footnotesize
\setlength{\tabcolsep}{4pt}
\begin{tabular}{lccccl}
\toprule
Region & $N_{\rm core}$ & $N>500$ & distance & topology & class \\
\midrule
Orion B   & 1870 & yes & yes & yes & primary \\
Perseus   & 816  & yes & yes & yes & primary \\
Aquila    & 749  & yes & yes & yes & primary \\
Taurus    & 536  & yes & yes & yes & primary \\
\midrule
Ophiuchus & 513  & yes & yes & \emph{no} & limited \\
Serpens   & ---  & yes & \emph{no} & yes & limited \\
TMC1      & 178  & \emph{no} & yes & yes & limited \\
CRA       & 239  & \emph{no} & yes & yes & limited \\
\bottomrule
\end{tabular}
\end{table}

\subsection{Measurement method}
\label{sec:methodology}

The analysis proceeds from the HGBS column-density maps and \textit{Gaia} DR3 distances through skeleton extraction, morphological bridging, core association within $0.06$\,pc of a spine, graph ordering along each component and adjacent-core separation. Filament skeletons are the published HGBS DisPerSE products \citep{Sousbie2011,Arzoumanian2011}, extracted with a $3\sigma$ persistence threshold and manual editing restricted to removing obvious artefacts. The effect of that editing on the spacing is small ($\lesssim5$ per cent), bounded by the junction-retention and association-threshold sensitivity tests; robustness to the persistence threshold itself is examined in Appendix~\ref{sec:persistence}. Because skeletons and core positions derive from distance-independent angular positions, the \textit{Gaia} revisions require no re-extraction. Excluding junctions changes the measured spacing by $<5$ per cent and varying the association threshold between one and $1.5$ widths by $<3$ per cent.

\smallskip\noindent\textbf{Projection, ties and junctions.} Catalogued sources are treated as points at their cataloged centroids. Source sizes, positional uncertainties and the covariance introduced by the common beam are \emph{not} propagated: the median catalogue positional uncertainty is a small fraction of a beam and is negligible against the spacings measured here, whereas the deblending limit, which is not a positional uncertainty but a resolution limit, is treated separately in Section~\ref{sec:pointprocess}. A source is assigned to the nearest skeleton pixel in projected distance. Where two branches are within one pixel of being equidistant the assignment is made to the branch of greater integrated column density along the ten pixels either side, and where a source lies within one association half-width of a junction it is assigned to one branch only, never counted twice; junction-adjacent sources are $3$--$7$ per cent of the associated sample and excluding them entirely changes the medians by less than $5$ per cent, as noted above. Spacings are computed within connected components and never across them.

Core-to-core spacings were also measured with the pairwise median, retained only as a comparability statistic with earlier work (Section~\ref{sec:statistics}).

\smallskip\noindent\textbf{Distances.} Physical spacing is angular separation times distance, so the $\lambda\propto d$ scaling is tautological and we do not present it as corroboration. The relevant check is that the residuals about that scaling show no trend with the absolute revision fraction, and the two most-revised clouds do not stand out. Two caveats remain. The YSO-based distances \citep{Zhang2023} are assumed to apply to the dense gas measured here; and the Aquila/Serpens Rift has substantial line-of-sight depth \citep{Zucker2020,Green2024}, so a single assigned distance may blend populations.

We quantified the blending with a geometric Monte-Carlo model in which two crests at different depths are projected onto a common spine, taking an intrinsic spacing of $0.20$\,pc at $d=436$\,pc. Each crest's offset is drawn from a Gaussian of dispersion $\sigma_z$, with the two separated by $2\sigma_z$, scanning $\sigma_z=5$--$30$\,pc. Superposition interleaves the two chains and \emph{shortens} the measured median, to $0.41$--$0.49$ of the true value across that range, whereas the distance spread alone contributes only $\pm1$ to $\pm7$ per cent. Blending biases the spacing towards shorter values rather than longer. The $\sim50$ per cent figure is a worst case in which two crests overlap completely. At the plausible end, $\sigma_z=5$--$10$\,pc with half the crest length overlapping, the shortening factor is $0.78$--$0.86$, which would move Aquila's $S_{\rm local}$ from $1.7$ to $2.0$--$2.2$ and $S_{\rm skel}$ from $3.4$ to $4.0$--$4.4$: still short locally and still below the reference band globally, so Aquila's classification survives a realistic correction. This is a geometric model with a free depth parameter, not a measurement; the stronger treatment, projecting the three-dimensional dust reconstructions along the actual sightline, is listed as a follow-up in Section~\ref{sec:limitations}. As a floor, retaining the pre-\textit{Gaia} distances lowers the primary-region characteristic to $S_{\rm local}\approx1.4$, so the sign of the result does not depend on the distance revision though its magnitude does.




\subsection{The choice of estimator}
\label{sec:statistics}

We use the along-crest adjacent-core spacing, writing $s_{\rm ACS}$ for its per-region median, and report the all-pairs median only for comparison with earlier work. It is not the Euclidean nearest-neighbour distance, which enters only in the Clark--Evans statistic of Section~\ref{sec:pointprocess}.

An all-pairs median, used as a local-spacing estimator, is dominated by the extent of the segment it is applied to. For $N$ points distributed over a segment of length $L$ it converges to $d^\star=L(1-1/\sqrt{2})\approx0.293\,L$, the median of the triangular pair-distance law, for a random and a perfectly periodic distribution alike (Appendix~\ref{sec:pwmedian}). It tracks that value for any filament with $N\gtrsim5$ cores and does not approach the local spacing, which for these clouds is one to two orders of magnitude smaller. The order statistic itself is elementary and long known; what matters here is its consequence for a statistic in use. Synthetic point processes confirm it: for cores placed by periodic, clustered, hierarchical and fractal processes with a fixed injected local scale, the adjacent-core median is flat with segment extent and recovers each process's local scale, whereas the all-pairs median rises as $0.29\,L$ for every one of them, independent of the injected scale.

Some methodological choices were fixed in advance and others compared afterwards, and the distinction matters. Fixed before any spacing was computed: the choice of adjacent-core over all-pairs separations, which follows from the argument above and not from the data; the primary/limited sample split, by the criteria of Section~\ref{sec:distance_uncertainties}; and the $0.06$\,pc association half-width, inherited from the HGBS filament width. Compared afterwards, and carrying no prior commitment: the width convention, the skeleton algorithm, and the injection--recovery correction branch. The systematics arising from those choices are collected in Appendix~\ref{sec:systematics}.

\subsection{Widths}
\label{sec:width}

Because $S_{\rm local}\propto1/W$, the width convention sets the normalisation. Two recur: the fixed-width convention ($W=0.10$\,pc for every region) and the region-Plummer convention, a beam-deconvolved Plummer FWHM fitted per region. We fit Plummer profiles perpendicular to the crest at every skeleton pixel, over a half-range of $0.4$\,pc with the outer $0.1$\,pc setting a local background, and deconvolve the $18.2''$ beam. Every region returns an index $p\simeq2.0$, the value found for HGBS filaments generally \citep{Arzoumanian2019}, and a width of $0.11$--$0.16$\,pc (Table~\ref{tab:widths}). The broader measured widths give a lower ratio, so the fixed $0.10$\,pc convention is the more conservative of the two and we quote results in it throughout, giving the region-Plummer values where the distinction matters.

\subsection{Validation}
\label{sec:injection_recovery}

Synthetic cores with periodic spacings of $0.15$--$0.35$\,pc were placed on skeletons matched to the HGBS properties and passed through the identical pipeline. Recovery is biased low by $11$--$29$ per cent, because bridging and association occasionally split or merge a pair, and the correction is applied as a per-region multiplicative factor. The size of that correction depends on the assumed point process: repeating the injection with a clustered process matched to the observed coefficient of variation gives a larger correction ($\div1.64$). That branch derives the correction by assuming the clustering it is meant to test, so we carry the periodic-injection value and quote the clustered branch only as a lower bound.

\section{Results}
\label{sec:results}

The four clouds are the independent statistical units throughout; the nesting of cores within clouds is propagated in the hierarchical bootstrap of Appendix~\ref{sec:systematics}, which resamples clouds, then skeleton components within clouds, then gaps within components.

\subsection{Two statistics, two answers}
\label{sec:twostats}

Table~\ref{tab:principal} gives the three numbers per cloud. Taking the unconditional statistic first, $S_{\rm skel}=3.4$, $3.4$, $4.5$ and $5.8$ for Orion~B, Aquila, Taurus and Perseus. Against the equilibrium-cylinder reference band of $5$--$6$, Perseus is consistent with it, Taurus lies just below, and Orion~B and Aquila fall below it by less than a factor of two. That proximity is not convention-invariant. It holds under the adopted persistent-crest network, the $0.06$\,pc association half-width, the all-core catalogue and the fixed-width normalisation; extracting the network as a mask medial axis instead moves $S_{\rm skel}$ to $9.1$--$16.7$, widening the association half-width to $0.30$\,pc moves it to $2.0$--$4.2$, and restricting the catalogue to prestellar cores moves it to $4.8$--$10.3$. The correct summary is therefore that $L_{\rm skel}/N$ lands near the equilibrium-cylinder reference for one defensible set of choices and not for others, and that its convention envelope spans a factor of eight.

Conditionally, $S_{\rm local}=1.8$, $1.7$, $3.0$ and $1.8$ in the same order, a factor of $\approx2$--$3.5$ below the reference band in every cloud and under every width and skeleton convention we tried (Appendix~\ref{sec:decoder}). Figure~\ref{fig:spacing} shows the per-region values in the two width conventions.

The factor of about two between the two statistics has two contributions, and neither is small-scale clustering. The median-to-mean rise of the gap distribution ($\approx30$--$120$ per cent) is the generic skew of a broad distribution, a pure exponential already giving $1.44$; the further rise to $L/N$ is dilution by extensive core-free filament. A local-width normalisation supports the same reading: fitting the beam-deconvolved transverse width at the midpoint of every adjacent-core interval ($146$--$545$ reliable fits per cloud) gives median $\Delta s/W_{\rm local}=0.61$, $0.91$, $1.05$ and $0.99$ for Taurus, Orion~B, Perseus and Aquila, all below the reference band, with only a weak spacing--width correlation ($r=-0.06$ to $0.31$).

Neither statistic is a measurement of a fragmentation wavelength. $S_{\rm skel}$ is the closer analogue of a one-fragment-per-wavelength model, but the equivalence holds only if the measured length corresponds to filament that was both fragmentation-eligible and velocity-coherent, which we cannot establish (Sections~\ref{sec:ontology} and \ref{sec:discussion}).

\subsection{Cores occupy a small part of the crest}
\label{sec:occupancy}


The gap between $S_{\rm local}$ and $S_{\rm skel}$ implies that cores occupy preferential stretches, and that can be measured. Each associated core is given an occupancy interval of $\pm0.06$\,pc along its spine; overlapping intervals are merged, and the occupancy fraction $f_{\rm occ}$ is the fraction of skeleton length that is core-bearing. We also count cores in $1$\,pc cells and form the Fano factor $F=\sigma^2/\mu$ of those counts (Table~\ref{tab:intermittency}, Fig.~\ref{fig:intermittency}).

On this operational definition $6$--$27$ per cent of the skeleton is core-bearing, so most filament length is empty at HGBS sensitivity. That number is tied to the adopted interval and is not an independently measured filling factor; the scale-independent evidence is the Fano behaviour and the distribution of core-free run lengths, below. The Fano factor exceeds unity everywhere ($1.6$, $2.8$, $5.7$, $25.2$ for Aquila, Orion~B, Perseus and Taurus). In Orion~B, $92$\,pc of the $337$\,pc skeleton is occupied, $155$\,pc lies in core-free runs inside core-bearing components, and $\approx90$\,pc is in components with no catalogued core at all.

Because $f_{\rm occ}$ grows with the assumed interval we report it against $\ell$ (Fig.~\ref{fig:scaletests}a): Taurus rises only from $0.052$ to $0.078$ and Perseus from $0.078$ to $0.202$ over a tenfold increase, so widening the interval adds almost nothing, whereas Orion~B and Aquila climb to $\approx0.6$. The Fano factor is the more robust statistic because it needs no occupancy convention: $F(\ell)>1$ across $\ell=0.25$--$4$\,pc in Taurus ($9.4$--$25$), Perseus ($2.3$--$8.4$) and Orion~B ($1.2$--$4.9$), rising monotonically with $\ell$. Aquila is the exception, consistent with Poisson below $1$\,pc.

\smallskip\noindent\textbf{Is this more than the filament's own structure implies?} A Fano factor above unity compares the cores with a \emph{homogeneous} Poisson process, and a filament is not homogeneous. The natural refinement is a null whose intensity is conditioned on the local column density. We construct one in Appendix~\ref{sec:nullspec} and show there that it cannot be calibrated: the cores raise the column density of their own surroundings, so the conditioning variable contains part of the phenomenon under test, and populations built with no clustering beyond the column-density structure are flagged as over-dispersed most of the time. We accordingly make no quantitative statement about whether the inhomogeneity exceeds that implied by the filament's own structure. Settling that question needs a tracer independent of the dust column, which means velocity-resolved dense-gas spectroscopy.

Two length scales coexist. On sub-parsec scales the gap distribution shows a deficit of the smallest separations (Section~\ref{sec:pointprocess}); on parsec scales the counts are strongly over-dispersed. Cores form in groups with short internal separations, and the groups are separated by long stretches that have produced none. This accounts for the factor of about two between $S_{\rm local}$ and $S_{\rm skel}$ without appealing to small-scale clustering, and it is the central observational result of this paper: what a fragmentation model must reproduce is a spatial point process, not a single spacing.






\begin{table}
\caption{Principal results, one row per cloud. $S_{\rm local}=s_{\rm ACS,med}/W$ is the conditional adjacent-core spacing, measured only within core-bearing stretches; $S_{\rm skel}=L_{\rm skel}/(NW)$ is the unconditional filament length per core; $S_{\rm elig}=(L_{\rm elig}/N_{\rm elig})/W$ restricts the latter to filament above the critical line mass, selected on a core-masked column-density reconstruction (Section~\ref{sec:eligible}), with $f_{\rm elig}$ the eligible length fraction and $N_{\rm elig}$ the cores on it; $f_{\rm occ}$ is the $0.06$\,pc occupancy fraction. All are normalised by the fixed $W=0.10$\,pc width so the columns are directly comparable, and all use the persistent-crest (DisPerSE) skeleton at a $0.06$\,pc association half-width. Brackets on $S_{\rm local}$ are $95$ per cent moving-block bootstrap intervals on the cloud median (blocks of five contiguous gaps, $4000$ resamples). They describe sampling uncertainty within a cloud, and are an order of magnitude smaller than the structural-convention envelope of Fig.~\ref{fig:envelope}, which is the term that limits the comparison. The equilibrium-cylinder reference is $5$--$6$. The simulated upper bound is common to all rows, $S_{\rm sim}^{\rm upper}=\lambda_{\rm turn}/W\lesssim1.0$--$1.3$ (Paper~II). $^{\dagger}$~The $S_{\rm elig}$ and $f_{\rm elig}$ columns are an \emph{exploratory diagnostic}, not a measurement: the eligibility criterion is a column-density proxy for a line mass and is partly determined by the cores it is used to interpret (Section~\ref{sec:eligible}); they should not be read on the same footing as $S_{\rm local}$ and $S_{\rm skel}$. The Perseus entry is bracketed because only $57$ cores remain on eligible filament after masking and the value depends on the masking radius (Table~\ref{tab:maskrad}); no value is quoted for Taurus, where $13$ cores remain and the statistic does not stabilise.}
\label{tab:principal}
\footnotesize
\setlength{\tabcolsep}{4pt}
\begin{tabular}{lccccc}
\toprule
Cloud & $S_{\rm local}$ [95\%] & $S_{\rm skel}$ & \emph{$S_{\rm elig}$}$^{\dagger}$ & \emph{$f_{\rm elig}$}$^{\dagger}$ [$N_{\rm elig}$] & $f_{\rm occ}$ \\
\midrule
Orion~B & 1.8 [1.6, 1.9] & 3.4 & \emph{2.0}$^{\dagger}$ & \emph{0.16 [268]} & 0.27 \\
Aquila  & 1.7 [1.5, 1.9] & 3.4 & \emph{3.2}$^{\dagger}$ & \emph{0.90 [303]} & 0.27 \\
Perseus & 3.0 [2.6, 3.6] & 5.8 & \emph{(4.1)}$^{\dagger}$ & \emph{0.08 [57]} & 0.12 \\
Taurus  & 1.8 [1.5, 2.2] & 4.5 & ---    & --- [13]   & 0.06 \\
\bottomrule
\end{tabular}
\end{table}

\begin{figure*}
\includegraphics[width=0.85\columnwidth]{figures/figure1_spacing_comparison.pdf}
\caption{The conditional local descriptor $S_{\rm local}=s_{\rm ACS,med}/W_{\rm fil}$ per primary region. This is \emph{not} the quantity the equilibrium-cylinder calculation gives (that is $S_{\rm skel}$; Fig.~\ref{fig:intermittency}b) and the comparison with the reference band below should be read diagnostically. It is the \emph{conditional} descriptor measured within core-bearing stretches; the unconditional length per core $L/N$, whose reciprocal is the line density predicted by fragmentation theory, is shown in Fig.~\ref{fig:intermittency}, in the two DisPerSE width conventions: directly-measured region-Plummer widths (blue circles, preferred; per-region raw $1.6$, $1.2$, $2.0$, $1.0$ for Taurus, Orion~B, Perseus and Aquila; the dashed line is their median, $1.4$) and the fixed $0.10$\,pc width (orange squares; $1.8$, $1.8$, $3.0$, $1.7$; median $1.8$), both shown \emph{raw} (before the injection--recovery correction; Section~\ref{sec:injection_recovery}). The red band is the equilibrium-cylinder reference wavelength $\lambda_{\rm m}\approx5$--$6\times$FWHM, exact for the $p=4$ Ostriker equilibrium; the conversion to FWHM is not invariant for the observed $p\simeq2$ profiles (\citealt{Fischera2012}; $=4\times$ the flat diameter, \citealt{Inutsuka1992}) and the dotted line the classical \emph{critical} wavelength $\lambda_{\rm cr}\approx2.6\times$FWHM; every region lies well below both. The grey band marks the FilFinder \citep{Koch2015} skeleton range ($\approx0.5$--$1.5$), and the vertical grey bars give the per-region min--max convention envelope, spanning the width conventions and the skeleton-algorithm factor. These are descriptive envelopes rather than confidence intervals (Section~\ref{sec:systematics}).}
\label{fig:spacing}
\end{figure*}

\begin{table}
\caption{Spatial-inhomogeneity diagnostics. $L_{\rm tot}$ is the total skeleton length, $L_{\rm core}$ the core-bearing length (union of $\pm0.06$\,pc intervals about each associated core), and $f_{\rm occ}$ is their ratio. $f_{\rm assoc}$ is the fraction of catalogue cores associated with the skeleton, and $F$ the Fano factor of counts in $1$\,pc cells. Both association half-widths are listed. }
\label{tab:intermittency}
\footnotesize
\setlength{\tabcolsep}{3pt}
\begin{tabular}{lccccccccc}
\toprule
Region & half & $L_{\rm tot}$ & $L_{\rm core}$ & $f_{\rm occ}$ & $f_{\rm assoc}$ & $\frac{L/N}{W}$ & $\frac{L_{\rm core}/N}{W}$ & $F$ & $G$ \\
       & (pc) & (pc) & (pc) & & & & & & \\
\midrule
Orion B & 0.06 & 337 & 92 & 0.27 & 0.53 & 3.4 & 0.9 & 2.8 & 0.49 \\
        & 0.30 & 337 & 117 & 0.35 & 0.85 & 2.1 & 0.7 & 4.5 & 0.50 \\
Aquila  & 0.06 & 106 & 29 & 0.27 & 0.41 & 3.4 & 0.9 & 1.6 & 0.40 \\
        & 0.30 & 106 & 38 & 0.36 & 0.70 & 2.0 & 0.7 & 2.1 & 0.36 \\
Perseus & 0.06 & 290 & 35 & 0.12 & 0.61 & 5.8 & 0.7 & 5.7 & 0.43 \\
        & 0.30 & 290 & 40 & 0.14 & 0.85 & 4.2 & 0.6 & 10.5 & 0.47 \\
Taurus  & 0.06 & 210 & 13 & 0.06 & 0.86 & 4.5 & 0.3 & 25.2 & 0.57 \\
        & 0.30 & 210 & 14 & 0.07 & 0.97 & 4.0 & 0.3 & 27.3 & 0.57 \\
\bottomrule
\end{tabular}
\end{table}

\begin{figure}
\centering
\includegraphics[width=0.9\columnwidth]{figures/fig_intermittency.pdf}
\caption{Spatial inhomogeneity of core occurrence along the filaments. \emph{(a)} The $0.06$\,pc occupancy fraction (bars, left axis) is $6$--$27$\,per cent, while the Fano factor of counts in $1$\,pc cells (red, right axis, logarithmic) exceeds the Poisson value of unity in every cloud. \emph{(b)} The three normalised spacing statistics: the local adjacent-core spacing is well below the reference band, the global $(L/N)/W$ approaches it, and the core-bearing $(L_{\rm core}/N)/W$ lies below both because the occupied stretches are densely populated. The separation between the blue and green curves measures how unevenly the cores occupy the crest.}
\label{fig:intermittency}
\end{figure}

\begin{figure*}
\centering
\includegraphics[width=0.94\textwidth]{figures/fig_scaletests.pdf}
\caption{Scale and convention tests of the inhomogeneity result. \emph{(a)} Occupancy fraction against the adopted half-width $\ell$ (dotted line: the $0.06$\,pc value quoted in the text). Taurus and Perseus saturate, so their low occupancy is not a consequence of the chosen $\ell$; Orion~B and Aquila rise steadily and are less strongly segregated. \emph{(b)} Fano factor against cell length. Over-dispersion ($F>1$, dashed) persists across $0.25$--$4$\,pc in Taurus, Perseus and Orion~B, but Aquila is consistent with Poisson below $1$\,pc. \emph{(c)} Convergence against association half-width: the conditional $s_{\rm ACS,med}/W$ (solid) plateaus, whereas the unconditional $(L/N)/W$ (dashed) falls as more cores are admitted, approaching the reference band (shaded) only in Perseus and Taurus.}
\label{fig:scaletests}
\end{figure*}

\subsection{The point process}
\label{sec:pointprocess}

Classical fragmentation predicts a quasi-periodic chain; fibre projection or a turbulent cascade predicts a broad distribution with no comb. Ordering associated cores by arclength and measuring adjacent gaps gives a coefficient of variation of order unity (Taurus $1.36$, Orion~B $0.83$, Perseus $0.96$, Aquila $0.65$; Fig.~\ref{fig:clustering}).

Goodness-of-fit tests were calibrated by parametric bootstrap ($N_{\rm boot}=10^5$). The result is two-sided: the adopted periodic-comb-plus-jitter model is rejected ($p_{\rm boot}<10^{-5}$ in every cloud, this being the floor set by the number of resamples rather than a measured value) \emph{and} so is the homogeneous Poisson (pure exponential) model ($p_{\rm boot}<10^{-5}$ in three clouds, $1.7\times10^{-3}$ in Perseus). Comparing five renewal models by the corrected Akaike criterion gives $\Delta$AICc$=37$--$246$ against Poisson and $81$--$742$ against the comb in every cloud. The shifted-exponential hard-core model is not uniquely preferred once broader alternatives are admitted, a positively-skewed lognormal fitting best in three clouds, so the data favour a broad, skewed distribution truncated at small separations, not a physical hard core.

These are, strictly, rejections of two specific nulls. Because a real fragmenting filament would not produce an ideal comb, we also tested a deliberately imperfect one: a chain whose local wavelength varies slowly along the crest, with positional jitter and a random fraction of fragments removed. Two versions were tried, a mild case ($\pm25$ per cent wavelength variation, $25$ per cent jitter, $15$ per cent dropout) and a strong case ($\pm50$ per cent, $25$ per cent, $35$ per cent). Both are rejected in all four clouds at the bootstrap floor ($p_{\rm boot}=3\times10^{-3}$, set by $300$ resamples), with Kolmogorov--Smirnov distances of $0.13$--$0.23$. The observed gap distributions are therefore inconsistent not only with an idealised comb but with a physically degraded one, which is a stronger statement than the AICc comparison alone supports. What remains unexcluded is a chain whose wavelength varies systematically with local filament properties rather than randomly, which would need those properties measured per stretch.

The clouds fail the exponential for different reasons and should not be summarised together (Table~\ref{tab:percloud}). Writing $R$ for the observed median gap divided by that expected for an exponential of the same mean, only Orion~B ($1.03$) and Aquila ($1.14$) show a deficit of the smallest separations; Perseus ($0.92$) and Taurus ($0.64$) show an excess.

\smallskip\noindent\textbf{The small-separation deficit.} The gap distributions are truncated below $0.011$--$0.044$\,pc. This range lies entirely inside the extraction floor: the $18.2''$ beam corresponds to $0.012$--$0.039$\,pc across our distance range and \textsc{getsources} \citep{Menshchikov2012} does not reliably separate sources closer than about one beam. The truncation is therefore \emph{consistent with} the extraction limit, and we attach no physical interpretation to it; we do not claim to have demonstrated that it is caused entirely by extraction, which would require the completeness experiment described below. It also bounds the interpretation of $S_{\rm local}$ itself: pairs closer than about one beam cannot be deblended, so any sub-beam hierarchical grouping, for instance fragments within a single velocity-coherent fibre, is invisible to this pipeline. The measured $S_{\rm local}$ is consequently an \emph{upper limit} on the true local separation of the gas fragments, not a measurement of it. Establishing the functional form of the gap distribution below the single-dish resolution requires interferometric imaging (ALMA or NOEMA) of the same crests. Deciding whether any physical suppression is also present requires the pair completeness function $C(s,F_1/F_2,N_{\rm bg})$, the probability that two sources of separation $s$, flux ratio $F_1/F_2$ and local background $N_{\rm bg}$ are recovered as two sources rather than one or none. It must be measured, not assumed, by injecting synthetic pairs spanning $s=0.3$--$3\,\theta_{\rm beam}d$, $F_1/F_2=1$--$10$ and the observed range of $N_{\rm bg}$ into the maps and re-extracting them through the unmodified HGBS pipeline. The observed gap density is then the true density multiplied by $C$, and only the deconvolved form is interpretable. That experiment is not attempted here, and in its absence the shape of the gap distribution below about $1.5\,\theta_{\rm beam}d$ is unconstrained: our measurements bound it from above but do not determine it.

\smallskip\noindent\textbf{Gap independence.} A runs test about the median gap shows significant non-independence in Perseus ($z=-4.2$) and Taurus ($z=-5.4$), the same core-poor/core-rich alternation documented in Section~\ref{sec:occupancy}. We therefore recomputed $\Delta$AICc with a moving-block bootstrap resampling contiguous runs of five adjacent gaps drawn from within single components ($400$ resamples per cloud). The effective sample size falls to $122$ of $440$ gaps in Taurus and $344$ of $438$ in Perseus, but the rejections survive: the $5$th percentile of the bootstrap $\Delta$AICc is $415$, $45$, $249$ and $587$ against the periodic model and $190$, $71$, $31$ and $55$ against Poisson. What does not survive is discrimination among the broad models, the shifted-exponential lying inside the lognormal's bootstrap interval in three clouds.

\smallskip\noindent\textbf{No classical-scale periodicity is detected.} The along-skeleton pair-correlation $g(s)$ (Fig.~\ref{fig:paircorr}) shows no feature at $5$--$6\,W$. To decide whether that absence is significant we generated a $95$ per cent envelope from the column-density-conditioned null of Appendix~\ref{sec:nullspec}: at the classical scale the observed $g$ lies \emph{inside} the envelope in every cloud (Orion~B $1.02$ in $[0.84,1.26]$; Aquila $0.89$ in $[0.58,1.52]$; Perseus $1.02$ in $[0.91,1.14]$; Taurus $1.00$ in $[0.91,1.10]$). This is a non-detection rather than a demonstration of absence. Two limitations bound what it can mean. First, the envelope is generated from the conditional null of Appendix~\ref{sec:nullspec}, which we show cannot be calibrated, so the quoted interval is an indicative spread and not a confidence region; the corresponding statement against a homogeneous null, which \emph{is} calibrated, is weaker, because a homogeneous expectation is not the right comparison where the source intensity varies strongly along the crest with column density, line mass and sensitivity. Second, absence of a feature relative to either null is not the same as absence of a periodic component after conditioning on the large-scale intensity: a modulation of a few per cent superposed on an intensity field that itself varies by an order of magnitude would not be detected here. What the calibrated tests of this section do establish is that the gap sequence is not drawn from a periodic comb, a deliberately imperfect comb or a homogeneous Poisson process; what they do not establish is an upper limit on a periodic component of an \emph{inhomogeneous} process. Setting that limit requires estimating a smoothly varying intensity along each component from an external predictor, simulating on the observed network topology with its selection function, and fitting a periodic modulation on top; velocity-resolved data are needed to supply a predictor that is not the dust column itself.








\begin{figure}
\centering
\includegraphics[width=0.9\columnwidth]{figures/fig_clustering.pdf}
\caption{Adjacent-core gap distributions along the skeletons of the four primary regions (histograms), with exponential/Poisson (solid) and regular/periodic (dashed) reference curves; The shaded region at small separations spans the \textit{Herschel} beam and \textsc{getsources} deblending scale for each cloud; \emph{the distribution below approximately one beam is not interpreted astrophysically}. Each panel gives the coefficient of variation CoV, the 1-D Clark--Evans ratio $R$, and the exponential goodness-of-fit $p$-value \emph{calibrated by parametric bootstrap} ($p_{\rm boot}$; the null distribution of the KS statistic is built by resampling from the fitted exponential, removing the estimated-parameter bias of the standard test). The result is two-sided: the regular/periodic model is rejected ($p_{\rm boot}<10^{-5}$, the parametric-bootstrap floor, $1/N_{\rm boot}$ with $N_{\rm boot}=10^5$) in every region, \emph{and} the pure-exponential (Poisson) model is also rejected ($p_{\rm boot}<10^{-5}$ (Perseus $1.7\times10^{-3}$) in all four), so in each of the four clouds the cores are neither periodic nor pure-Poisson; consistent with broad, cloud-dependent gap distributions, affected by extraction incompleteness at the smallest separations and over-dispersed on parsec scales. The direction of the small-gap departure is not common to the sample: $R>1$ in Orion~B and Aquila indicates a deficit of small gaps relative to an exponential, whereas $R<1$ in Perseus and Taurus indicates an excess (Table~\ref{tab:percloud}).}
\label{fig:clustering}
\end{figure}

\begin{figure}
\centering
\includegraphics[width=0.9\columnwidth]{figures/fig_paircorr.pdf}
\caption{Along-skeleton pair-correlation function $g(s)$ for the four primary clouds (uniform-random expectation on the matched components divided out; $g=1$ is no correlation). Dotted vertical lines mark the classical scale $5.5\,W_{\rm fil}$ per region. No cloud shows a peak at the classical scale ($g\approx1$ there, slightly below unity in Orion~B and Aquila); the only departures are a mild small-separation excess (Taurus) and weak region-envelope features near $1$\,pc. No feature is detected at the classical scale. This does not exclude a weak periodic modulation superposed on an inhomogeneous intensity field, which these data cannot bound (Section~\ref{sec:pointprocess}).}
\label{fig:paircorr}
\end{figure}

\begin{table}
\caption{Cloud-by-cloud summary of the point-process diagnostics, to avoid summarising four different behaviours as one. $R$ is the observed median gap divided by the median expected for an exponential of the same mean ($R>1$ deficit of small gaps, $R<1$ excess); $F_{\rm obs}$ is the observed Fano factor of counts in $1$\,pc cells and median $F_{\rm null}$ that of the inhomogeneous, column-density-conditioned null of Appendix~\ref{sec:nullspec}. The last two columns are \emph{not} probabilities and must not be read as such. They are null-exceedance fractions, $E_{\rm fit}$ using an intensity exponent fitted to the cloud's own cores and $E_{\rm cv}$ one predicted from the other three under leave-one-cloud-out cross-validation, and the null generating them has an uncontrolled false-positive rate, so no significance attaches to any entry (Section~\ref{sec:occupancy}, Appendix~\ref{sec:nullspec}). They are tabulated only as descriptive discrepancy ranks. The calibrated tests in this paper are those of Section~\ref{sec:pointprocess} against the periodic and Poisson nulls; these are not among them.}
\label{tab:percloud}
\footnotesize
\setlength{\tabcolsep}{3.5pt}
\begin{tabular}{lcccccc}
\toprule
Region & CoV & $R$ & $F_{\rm obs}$ & median $F_{\rm null}$ & $E_{\rm fit}^{\ddagger}$ & $E_{\rm cv}^{\ddagger}$ \\
\midrule
Orion B & 0.83 & 1.03 & 2.8  & 3.6  & 0.98 & 1.00 \\
Aquila  & 0.65 & 1.14 & 1.6  & 2.1  & 0.88 & 0.27 \\
Perseus & 0.96 & 0.92 & 5.7  & 7.1  & 0.96 & 0.995 \\
Taurus  & 1.36 & 0.64 & 25.2 & 13.4 & $2\times10^{-4}$ & $0.005$ \\
\bottomrule
\end{tabular}
\\[2pt]
\footnotesize $^{\ddagger}$\,\emph{Not p-values.} $E_{\rm fit}$ and $E_{\rm cv}$ are uncalibrated null-exceedance fractions from a conditional null whose own false-positive rate is uncontrolled. No significance attaches to them and they should not be quoted as probabilities.
\end{table}


\section{What limits the answer}
\label{sec:ontology}

\subsection{The filament is not uniquely defined}


A skeleton defines $L$ itself, so changing the extraction algorithm alters the spacing, the length, the core association and the topology together. DisPerSE \citep{Sousbie2011} and FilFinder \citep{Koch2015} are not two noisy estimates of one observable but partly different definitions of the object, and we therefore recomputed \emph{all} the principal statistics under both (Table~\ref{tab:ontology}).

The FilFinder networks are shorter ($67$--$121$\,pc against $106$--$337$\,pc) and associate far fewer catalogue cores, so both statistics move, but they move in \emph{opposite} directions: $S_{\rm local}$ becomes shorter still ($0.54$--$1.54$ against $1.7$--$3.0$) while $S_{\rm skel}$ becomes much larger ($9.1$--$16.7$ against $3.4$--$5.8$). The local/skeleton contrast survives the change of structural definition and in fact widens, from a factor $\approx2$ to $\approx9$--$17$, and the occupancy fraction stays low under both. We regard this as one of the more secure results here, precisely because the two algorithms disagree so strongly about what a filament is.

At a high-column hub, FilFinder recovers the main ridge but adds long, almost straight spokes that leave the emission entirely (Fig.~\ref{fig:skeleton_compare}). These are most likely medial-axis artefacts of a large connected mask rather than resolved sub-filaments, though we have not quantified what fraction of the total length behaves this way and column density alone cannot decide: position--position--velocity data are needed to establish which network corresponds to a coherent physical object. We therefore base our quantitative conclusions on the persistent-crest network and carry FilFinder as an ontology test.

\smallskip\noindent\textbf{Coverage bias.} Between $41$ and $86$ per cent of catalogue cores associate with the skeleton (Table~\ref{tab:nnpipeline}), and the rejected cores have projected separations $1.5$--$3.0$ times \emph{larger}, so coverage selection biases the adjacent spacing short. Relaxing the association half-width from $0.06$ to $0.30$\,pc raises coverage to $70$--$97$ per cent and lengthens the median adjacent spacing by only $3$--$17$ per cent (Fig.~\ref{fig:scaletests}c); $S_{\rm local}$ plateaus and the bias is bounded at $\lesssim17$ per cent.

The unconditional statistic behaves differently, and this is a genuine limitation. Because adding cores at fixed $L$ necessarily raises the line density, $S_{\rm skel}$ falls monotonically with association radius, from $3.6$ to $2.1$ in Orion~B, $3.6$ to $2.0$ in Aquila, $6.6$ to $4.2$ in Perseus and $5.1$ to $4.0$ in Taurus. The near-classical value is robust in Perseus and Taurus but in Orion~B and Aquila holds only near the adopted radius. The local/global contrast persists in all four clouds; its size in the two low-coverage clouds depends on how generously cores are assigned.






\begin{table}
\caption{Principal statistics under both filament definitions, at a $0.06$\,pc association half-width and a $0.06$\,pc occupancy half-width. $F(1\,{\rm pc})$ is omitted where the FilFinder components are mostly shorter than the cell. The local and global statistics move in opposite senses between the two algorithms, so their contrast is not an artefact of either ontology.}
\label{tab:ontology}
\footnotesize
\setlength{\tabcolsep}{3pt}
\begin{tabular}{llccccc}
\toprule
Region & Skeleton & $L$ & $N_{\rm assoc}$ & $(L/N)/W$ & $s_{\rm ACS,med}/W$ & $f_{\rm occ}$ \\
 & & (pc) & & & & \\
\midrule
Orion B & DisPerSE  & 337 & 995 & 3.4 & 1.77 & 0.27 \\
        & FilFinder & 121 &  73 & 16.6 & 1.54 & 0.06 \\
Aquila  & DisPerSE  & 106 & 308 & 3.4 & 1.66 & 0.27 \\
        & FilFinder &  67 &  40 & 16.7 & 0.97 & 0.05 \\
Perseus & DisPerSE  & 290 & 501 & 5.8 & 2.97 & 0.12 \\
        & FilFinder &  78 &  85 &  9.1 & 1.05 & 0.09 \\
Taurus  & DisPerSE  & 209 & 463 & 4.5 & 1.81 & 0.06 \\
        & FilFinder & 111 & 120 &  9.3 & 0.54 & 0.05 \\
\bottomrule
\end{tabular}
\end{table}

\begin{figure}
\centering
\includegraphics[width=0.92\columnwidth]{figures/fig_skeleton_compare.pdf}
\caption{The dominant observational systematic, shown directly: the DisPerSE (blue) and FilFinder (red) skeletons extracted from the \emph{same} \textit{Herschel} column-density map of Orion~B (greyscale, $\log N_{\rm H_2}$; a representative sub-region). The FilFinder skeleton is produced by \textsc{FilFinder2D} with the cloud distance set per region and an \emph{externally supplied} mask (\texttt{use\_existing\_mask=True}) followed by \texttt{medskel()}; the mask is a simple column-density cut at the $92$nd percentile of the finite pixels, which corresponds to $N_{\rm H_2}=2.5$, $3.8$, $3.1$ and $2.9\times10^{21}$\,cm$^{-2}$ in Orion~B, Aquila, Perseus and Taurus respectively. Because the mask is supplied externally, \textsc{FilFinder}'s own adaptive and branch thresholds are not invoked. A more permissive $70$th-percentile mask is also computed and is what sets the upper end of the FilFinder range. The DisPerSE skeletons are the published HGBS products, extracted at a $3\sigma$ persistence cut. The two algorithms trace substantially different filament networks, FilFinder's medial axis of a column-density mask is denser and more branched and follows more low-contrast material, whereas DisPerSE follows the persistent crests, which is why the inferred $S_{\rm local}$ differs by a factor of $\sim$2 between them (Section~\ref{sec:systematics}). This is intrinsic to how each algorithm defines a ``filament'' and is the largest single uncertainty on the \emph{magnitude} of the spacing, though the direction of the offset is the same for both.}
\label{fig:skeleton_compare}
\end{figure}

\subsection{Why this is a limit on the comparison, not an error bar}
\label{sec:whatfilament}

The factor-of-two DisPerSE-versus-FilFinder difference (Fig.~\ref{fig:skeleton_compare}) is not measurement noise: the two algorithms \emph{define different objects} (persistent crests versus the medial axis of a masked region), so there is no single, algorithm-independent ``filament'' whose spacing they should reproduce. The \emph{sign} of the deficit is robust, every convention placing $S_{\rm local}$ below the reference band, but its \emph{magnitude} has no unique value: the same data yield $\approx0.7$, $\approx1.4$ or $\approx1.8$ under three defensible combinations of skeleton and width. Segmentation alone moves the answer by as much as the differences between the physical models compared in Paper~II, so it cannot discriminate between them (Fig.~\ref{fig:envelope}).



The deeper point was stated in Section~\ref{sec:intro}: neither construction guarantees that a single coherent cylinder exists along the traced path, and that cannot be established from column density. The comparison in this paper is a diagnostic and not a test, and the decisive follow-up is kinematic.


\begin{figure*}
\centering
\includegraphics[width=0.92\textwidth]{figures/fig_envelope.pdf}
\caption{Observed statistics and model scales, kept apart. \emph{(a)} The point-process quantities measured here: $S_{\rm local}$ under three structural and width conventions, $S_{\rm skel}$ under two skeleton definitions, and the exploratory $S_{\rm elig}$. Each bar spans the range across the four clouds. \emph{(b)} Scales predicted by models, with the observable each would predict indicated: a one-bead-per-wavelength calculation predicts a line density and so should be set against $S_{\rm skel}$, whereas the hierarchical-fibre construction predicts the within-group separation and is closer to $S_{\rm local}$. Neither correspondence is exact: $S_{\rm skel}$ is the closest available proxy for a line density only if the measured length is fragmentation-eligible and velocity-coherent, which is not established here. The simulated entry is an upper limit and is drawn as an arrow. The grey band is the equilibrium-cylinder reference $\lambda_{\rm m}/W_{\rm FWHM}\approx5$--$6$ for the projected $p=4$ profile, exact for that equilibrium and not invariant for the observed $p\simeq2$ profiles; the dotted line is $\lambda_{\rm cr}$. No quantitative identification between a quantity in (a) and one in (b) is intended: they are shown on a common dimensionless axis only to compare magnitudes. The horizontal axis is symmetric-log below $\lambda/W=1$.}
\label{fig:envelope}
\end{figure*}


\subsection{Coverage, population and the other systematics}
\label{sec:coverage}

\begin{table}
\caption{Headline statistics as a function of the core--skeleton association half-width, reported at five values rather than only at the adopted $0.06$\,pc. $f_{\rm assoc}$ is the fraction of catalogued sources assigned to the skeleton. $S_{\rm local}$ varies by at most $17$ per cent across the whole range in every cloud; $S_{\rm skel}$ falls monotonically, because widening the association adds cores at fixed skeleton length. The conditional statistic is robust to this choice and the unconditional one is not.}
\label{tab:radsweep}
\footnotesize
\setlength{\tabcolsep}{4pt}
\begin{tabular}{lccccc}
\toprule
Cloud & half (pc) & $f_{\rm assoc}$ & $S_{\rm local}$ & $S_{\rm skel}$ & $f_{\rm occ}$ \\
\midrule
Orion B & 0.04 & 0.51 & 1.86 & 3.56 & 0.26 \\
        & \textbf{0.06} & 0.53 & \textbf{1.77} & \textbf{3.39} & 0.27 \\
        & 0.10 & 0.60 & 1.70 & 3.03 & 0.29 \\
        & 0.20 & 0.77 & 1.83 & 2.34 & 0.33 \\
        & 0.30 & 0.85 & 1.96 & 2.12 & 0.35 \\
\midrule
Aquila  & 0.04 & 0.39 & 1.68 & 3.59 & 0.26 \\
        & \textbf{0.06} & 0.41 & \textbf{1.66} & \textbf{3.44} & 0.27 \\
        & 0.10 & 0.45 & 1.58 & 3.16 & 0.28 \\
        & 0.20 & 0.58 & 1.72 & 2.43 & 0.33 \\
        & 0.30 & 0.70 & 1.94 & 2.03 & 0.36 \\
\midrule
Perseus & 0.04 & 0.54 & 3.19 & 6.61 & 0.11 \\
        & \textbf{0.06} & 0.61 & \textbf{2.97} & \textbf{5.78} & 0.12 \\
        & 0.10 & 0.72 & 3.17 & 4.95 & 0.13 \\
        & 0.20 & 0.81 & 3.09 & 4.41 & 0.13 \\
        & 0.30 & 0.85 & 3.20 & 4.19 & 0.14 \\
\midrule
Taurus  & 0.04 & 0.77 & 1.76 & 5.07 & 0.06 \\
        & \textbf{0.06} & 0.86 & \textbf{1.81} & \textbf{4.52} & 0.06 \\
        & 0.10 & 0.92 & 1.84 & 4.23 & 0.07 \\
        & 0.20 & 0.96 & 1.82 & 4.09 & 0.07 \\
        & 0.30 & 0.97 & 1.86 & 4.04 & 0.07 \\
\bottomrule
\end{tabular}
\end{table}


The adjacent-core measurement thresholds the deposited DisPerSE skeleton, bridges fragmented segments by morphological closing and re-skeletonisation, associates each core to the nearest spine within $0.06$\,pc, orders the associated cores along each component by graph arclength (double-BFS diameter) and takes adjacent separations. An adjacent pair is defined within a single connected skeleton component: cores are projected onto the component's longest geodesic and ordered along it, so a pair straddling a junction is retained if both cores project onto the same geodesic and is otherwise not formed. Pairs are never formed between components. This is the convention that inflates the \emph{mean} adjacent spacing on heavily bridged skeletons, which is why we use the median. Table~\ref{tab:nnpipeline} lists the resulting counts.



The associated fraction ranges from $0.41$ to $0.86$, so in Orion~B and Aquila the spacing is measured on roughly half the catalogue. Catalogue incompleteness is itself spatially dependent, because source extraction is harder against a high and structured background and in crowded regions, so the densest stretches are the ones where cores are most likely to be missed or merged. Two distinct effects are involved and we keep them apart. The first is \emph{geometric association}: whether a catalogued source is assigned to the skeleton at all. This we can bound, because it is controlled by a parameter we can vary. Relaxing the association half-width to $0.30$\,pc raises the coverage to $70$--$97$ per cent while lengthening the median adjacent spacing by only $3$--$17$ per cent, and sweeping it over $0.04$--$0.08$\,pc changes the per-region median by $4$--$7$ per cent, so the geometric term alone is one-sided and bounded at $\lesssim17$ per cent. Table~\ref{tab:radsweep} reports the headline statistics at four half-widths rather than only at the adopted one. The second effect is \emph{source-detection completeness}, which is not bounded by any of these sweeps: extraction is hardest against a high, structured background and in crowded stretches, so cores are preferentially missed or merged exactly where the spacing is shortest. Its sign is the same, towards a short measured spacing, but its size is unknown without the map-level injection experiment of Section~\ref{sec:pointprocess}, and we make no quantitative claim about it. The $\lesssim17$ per cent figure therefore bounds the association step and not the total selection effect.

Restricting to the two clouds with the highest coverage, Perseus and Taurus, gives the same median as the full four-region sample. The dependence on the core population is examined in Section~\ref{sec:coverage}.



\begin{table}
\caption{Measurement counts per primary region: catalogue cores, cores associated to a skeleton within $0.06$\,pc, the associated fraction, the number of adjacent-core pairs entering the median, and the number of connected skeleton components retained, with the number that carry at least two associated cores and so contribute pairs in brackets. The clouds differ as much in topology as in core number: Taurus's $440$ pairs come from only $11$ components, Orion~B's $699$ from $203$.}
\label{tab:nnpipeline}
\footnotesize
\setlength{\tabcolsep}{5pt}
\begin{tabular}{lccccc}
\toprule
Region & catalogue & associated & $f_{\rm assoc}$ & pairs & components \\
\midrule
Orion B & 1870 & 995 & 0.53 & 699 & 378 (203) \\
Aquila  &  749 & 308 & 0.41 & 168 & 182 (77) \\
Perseus &  816 & 501 & 0.61 & 438 & 107 (32) \\
Taurus  &  536 & 463 & 0.86 & 440 & 104 (11) \\
\bottomrule
\end{tabular}
\end{table}


The HGBS catalogues contain unbound starless cores, prestellar cores and protostellar sources. These are not equivalent for the present purpose: protostellar objects have had time to accrete, migrate and interact, while unbound starless cores need not be fragmentation products at all. We repeated the principal statistics on three nested samples per cloud (Table~\ref{tab:popsplit}).

The conditional statistic is essentially invariant. $S_{\rm local}$ changes by at most $0.2$ between the all-core and prestellar-only samples in every cloud, and the coefficient of variation of the gaps is unchanged to within $0.1$. The unconditional statistic behaves in the opposite way, rising steeply as the sample is restricted, from $3.4$ to $5.9$ in Orion~B, $3.4$ to $4.8$ in Aquila and $5.8$ to $10.3$ in Perseus, simply because $L$ is fixed while $N$ falls.

That combination sharpens the central result rather than weakening it. The separation between the local and global statistics is not produced by the inclusion of evolved or unbound objects; restricting to the cores most likely to be genuine fragmentation products \emph{widens} it, and pushes $S_{\rm skel}$ from near-classical to super-classical in three of the four clouds. It also shows that $S_{\rm skel}$ is sensitive to catalogue selection in a way that $S_{\rm local}$ is not, which is a further reason to treat the comparison of $S_{\rm skel}$ with $\lambda_{\rm m}$ as a diagnostic and not a measurement. The Taurus prestellar sample contains $54$ catalogue cores, of which $53$ associate with a skeleton, so its $S_{\rm skel}$ is not meaningful and is quoted only for completeness.




The remaining systematic terms, the width and persistence conventions, the injection--recovery correction branch, deprojection, line-of-sight blending and the manual editing of the published skeletons, are collected in Appendix~\ref{sec:systematics}. None of them reverses a conclusion; they move numbers within the envelope of Fig.~\ref{fig:envelope}.

\subsection{An exploratory diagnostic: how much of the crest could fragment?}
\label{sec:eligible}


Identifying $L/N$ with the quantity a one-fragment-per-wavelength theory predicts assumes that the whole skeleton could have fragmented; a subcritical stretch has no axial instability to express. The estimate below is exploratory. The criterion is a column-density proxy for a line mass and is partly determined by the cores it is used to interpret, and none of the conclusions of this paper rests on it.

We define eligibility \emph{independently of where the cores are}, using column density. A filament of width $0.10$\,pc reaches $M_{\rm line,crit}=2c_s^2/G\simeq16\,M_\odot\,{\rm pc^{-1}}$ at $N_{\rm H_2}\simeq7.2\times10^{21}\,{\rm cm^{-2}}$ ($\mu=2.8$), and we take $L_{\rm elig}$ to be the skeleton length above that threshold. This is a proxy, not a line-mass measurement: the conversion depends on width, inclination, temperature, background subtraction and both non-thermal and magnetic support. The eligible fractions are $0.19$, $0.91$, $0.13$ and $0.07$ for Orion~B, Aquila, Perseus and Taurus, so Orion~B in particular is dominated by filament that could not have fragmented at all.

\smallskip\noindent\textbf{The criterion is partly endogenous.} A core raises the column density of its own surroundings, so the threshold can be crossed \emph{by the core it is meant to be independent of}. Between $25$ and $34$ per cent of skeleton pixels lie within one association half-width of a core. We therefore rebuilt the criterion on a core-masked background, replacing the column density within a masking radius of each associated core by interpolation along the spine from untainted pixels. On that background $S_{\rm elig}$ becomes $2.0$, $3.2$ and $4.1$ for Orion~B, Aquila and Perseus, against the raw $1.6$, $3.2$ and $2.0$, and $N_{\rm elig}$ falls from $199$ to $57$ in Perseus and from $186$ to $13$ in Taurus. Thirteen surviving cores are too few for a median, so no masked value is quoted for Taurus; the Perseus value is bracketed throughout for the same reason in weaker form.

The statistic is also sensitive to the masking radius (Appendix~\ref{sec:maskapp}). Only Orion~B and Aquila are comparatively stable against it and retain more than $230$ eligible cores; Perseus drifts over $3.6$--$4.5$ on $57$ cores, and Taurus does not stabilise at any radius. Nor has the masking step itself been validated by injection and recovery. The endogeneity is identified and mitigated, not removed. We therefore quote $S_{\rm elig}$ for Orion~B and Aquila only, bracket Perseus and omit Taurus, and we do not use any of them to support a comparison with the equilibrium-cylinder scale. The underlying difficulty is not one of estimator design: a dust column cannot separate the mass of the filament from the local gravitational enhancement produced by the cores embedded in it, and no masking prescription applied to a column-density map can. Deciding which stretches of crest were ever unstable requires a kinematic tracer.

\smallskip\noindent\textbf{What distinguishes a fertile stretch from a sterile one?} Supercriticality is necessary but not sufficient. In Aquila $90$ per cent of the crest is eligible yet only $24$ per cent carries a core. Local compression, longitudinal velocity gradients of $\sim1$\,km\,s$^{-1}$\,pc$^{-1}$ \citep{Kirk2013,Hacar2013}, and spatially varying non-thermal support are all plausible controls, and none can be tested with column density alone. The occupancy is thus reported as a measured property whose cause is undetermined.

\smallskip\noindent\textbf{A prescription for future use.} If this diagnostic is to be comparable between studies the masking radius must be tied to a physical scale. We recommend $r_{\rm mask}=\max(W_{\rm fil},\theta_{\rm beam}d)$, one filament FWHM subject to a floor of one beam, which for our sample is $0.11$--$0.16$\,pc; and that the statistic be reported with $N_{\rm elig}$ and not quoted where $N_{\rm elig}\lesssim100$. Separating fertile from sterile stretches is not a column-density problem at all, and the observation that would settle it is set out in Section~\ref{sec:settle}.

A physically motivated restriction to nominally unstable crest moves $S_{\rm local}$ and $S_{\rm skel}$ towards each other, but current column-density data do not permit a sufficiently independent determination of the unstable filament length.




\section{Discussion}
\label{sec:discussion}

\subsection{What the observations do and do not establish}

Three statements are supported by the measurements above, and they are the three numbers of Table~\ref{tab:principal}: the conditional adjacent-core spacing is several times below the equilibrium-cylinder reference in every cloud; the skeleton length per core lies near that reference in two clouds under the adopted conventions and not under others; and the cores occupy a small, clustered fraction of the crest, with a gap distribution inconsistent with both a periodic comb and a homogeneous Poisson process and with no feature at the reference scale.

Two statements are not supported. We have not measured a fragmentation wavelength: an evolved point process need not preserve a one-to-one mapping onto the modes that produced it, since subsequent accretion, merging, migration and incompleteness all break it. And we have not established that the traced skeletons are the objects to which a cylinder eigenmode applies.

\subsection{What could produce a short local spacing}

At least four causes can produce a short $S_{\rm local}$, and column density alone cannot rank them.

\begin{itemize}\setlength{\itemsep}{1pt}
\item A genuinely short fragmentation wavelength, from magnetic tension in a force-balanced helical field, or from fragmentation of a filament that is not in equilibrium. One theoretical assumption relevant to the latter, that an evolving filament must select the wavelength of its initial equilibrium, is tested in Paper~II; that paper supplies a counterexample to the assumption and does not supply an explanation of the observations.
\item Several velocity-coherent fibres superposed within one morphological filament, so that cores from different fibres interleave on a common projected crest.
\item Spatially varying conditions along the crest, so that the populated stretches are simply those that happened to be unstable.
\item Post-fragmentation evolution: accretion, migration and merging over several free-fall times \citep{Kirk2016,Mattern2018,Clarke2020}.
\end{itemize}

The second of these deserves a note because it is quantitatively sufficient on its own. If a filament contains $N$ fibres of width $W_{\rm fib}$, each fragmenting near its own equilibrium scale, and the cores are measured across the bundle, then in the interleaving limit
\begin{equation}
\frac{\lambda}{W_{\rm fil}}\approx\frac{(5\text{--}6)\,r}{N},\qquad r\equiv W_{\rm fib}/W_{\rm fil},
\label{eq:fourrn}
\end{equation}
which reproduces the observed values for $r/N\approx0.2$--$0.4$ on the adopted convention, and for $r/N\approx0.1$--$0.6$ across the full convention envelope. Equation~\ref{eq:fourrn} is a geometrical scaling in the interleaving limit and not a fragmentation model: real fibres need not be independent, of equal width or randomly phased, and may overlap in position and velocity, interact gravitationally, differ in line mass and evolve. What it shows is that fibre multiplicity \emph{can} generate apparently short filament-scale spacings without any change to the fragmentation physics, not that it quantitatively accounts for the observed values.

\subsection{What would settle it}
\label{sec:settle}

One measurement addresses the structural and the physical ambiguity together. Velocity-resolved dense-gas spectroscopy, N$_2$H$^+$(1--0) or NH$_3$(1,1) at $\lesssim0.05$\,pc resolution, separates coherent fibres from projected superpositions, gives the local non-thermal linewidth and hence the virial line mass per stretch, and supplies a conditioning variable for the clustering null that is not the same quantity as the dust column. Concretely we suggest: decompose the cube into velocity-coherent components by grouping adjacent pixels whose centroid velocities differ by less than the local linewidth; extract a skeleton from each component separately; retain a segment only where the gas along it is coherent in velocity over its length; and compute spacing statistics per component. Applied to a medial-axis network this removes the long straight spokes of Fig.~\ref{fig:skeleton_compare}, which traverse material at different velocities; applied to a persistent-crest network it should change little, which is itself the test of whether the two definitions converge once kinematics are imposed.

Until such data exist, we recommend that any quoted filament spacing state the extraction algorithm and its parameters, and give the value under at least two structural definitions.

\subsection{Limitations}
\label{sec:limitations}


\emph{Sample.} Four nearby, low-mass Gould Belt regions of broadly solar metallicity and similar radiation field are not a population, and whether the result persists in high-mass complexes, at lower metallicity or under stronger irradiation is untested. How much of the cloud-to-cloud spread could be finite-sample noise is nevertheless answerable: pooling all $1745$ adjacent-core gaps from the four clouds, drawing four samples of the observed per-cloud sizes ($699$, $168$, $438$, $440$) \emph{with} replacement, taking each sample's median and recording the range across the four, and repeating $10^{5}$ times, gives a median range in $S_{\rm local}$ of $0.24$ and a $95$th percentile of $0.48$, against an observed range of $1.31$; the observed range is exceeded in $5\times10^{-5}$ of realisations. The four measured clouds therefore differ by more than that particular pooled-gap resampling null expects; this is not a demonstration of cloud-to-cloud diversity across the Gould Belt population, which four clouds cannot support, nor does it identify which cloud property drives the differences.

\emph{Method.} Two validation tests are outstanding. The $\lesssim5$ per cent bound on manual skeleton editing is a proxy, set by the junction-retention and association-radius sensitivities rather than by a blind multi-editor re-skeletonisation; and the small-separation deficit cannot be separated into instrumental and physical parts without a map-level injection experiment. The core-masking step of Section~\ref{sec:eligible} is likewise unvalidated by injection and recovery. All code and catalogues are deposited so that these can be carried out.

\emph{Observations not available to us.} A per-segment line-mass split, which would test whether the most supercritical filaments show the shortest spacing, needs segment-level line masses at a fidelity the data do not support; and resolved velocity dispersions per stretch, which would say why some eligible filament forms cores and some does not, do not exist for these regions.

\smallskip\noindent\emph{Named follow-ups.} Four experiments would settle points left open here: projecting the three-dimensional dust reconstructions of \citet{Zucker2020} and \citet{Green2024} along the sightline through the Aquila/Serpens Rift, replacing the geometric blending model of Section~\ref{sec:methodology} by a count of physically distinct crests within one association length; injecting synthetic source pairs into the \textit{Herschel} maps and re-extracting them, to measure the close-pair completeness function; dense-gas position--position--velocity mapping to establish which skeleton corresponds to a velocity-coherent object; and, on the theoretical side, the calculations set out in Paper~II.

\smallskip\noindent\emph{A recommendation for observers.} Because the skeleton systematic dominates every comparison with theory, we recommend that skeletons be extracted to a fixed column-density contrast relative to the local background and not to an absolute level; that spacings be quoted for \emph{both} a persistent-crest and a medial-axis definition; and that any comparison with a one-dimensional cylinder model be based on velocity-coherent ridges identified in a dense-gas tracer instead of on a column-density mask.




\section{Conclusions}
\label{sec:conclusions}

We have measured core spacing along filament crests in four HGBS clouds with \textit{Gaia} DR3 distances, and asked whether it is consistent with the fastest-growing mode of an equilibrium isothermal cylinder.

\begin{enumerate}\setlength{\itemsep}{2pt}
\item \textbf{There is no single observational core spacing.} Adjacent cores within populated stretches are separated by $S_{\rm local}=1.7$--$3.0$ filament widths; the filament length per core over the whole crest is $S_{\rm skel}=3.4$--$5.8$. The two answer different questions and differ by about a factor of two.
\item \textbf{The proximity of the unconditional statistic to the classical scale is convention-dependent.} Under the adopted persistent-crest network, $0.06$\,pc association radius, all-source catalogue and fixed-width convention, $S_{\rm skel}$ lies near the equilibrium-cylinder reference in two clouds: consistent with it in Perseus, just below in Taurus, and below by less than a factor of two in Orion~B and Aquila. That proximity is not invariant. Other defensible skeleton, association and population choices move $S_{\rm skel}$ over a factor-of-eight envelope. It is therefore an operational diagnostic and not a test of the cylinder model.
\item \textbf{The difference between the two statistics is the result.} Cores occupy $6$--$27$ per cent of the crest on a $\pm0.06$\,pc definition, forming tight groups separated by long empty stretches. Core occurrence is spatially inhomogeneous, and what a fragmentation model must reproduce is a point process, not a spacing.
\item \textbf{The cores are neither a periodic comb nor a homogeneous Poisson process.} Both are rejected in all four clouds, as is a deliberately imperfect comb with a varying wavelength and dropouts, and the rejections survive a moving-block bootstrap. No feature is detected at the equilibrium-cylinder scale in the pair-correlation function. Whether the inhomogeneity exceeds that implied by the filament's own column density cannot be settled with these data.
\item \textbf{The inferred magnitude is dominated by the adopted structural definition.} Changing from a persistent-crest network to a mask medial axis moves $S_{\rm local}$ by a factor of two and $S_{\rm skel}$ by more, whereas persistence thresholding within one definition changes it by less than ten per cent. The sign of the offset is robust; its magnitude is not, and these data cannot separate the physical contribution from the definitional one. The limiting uncertainty on a dimensionless core spacing is the structural definition of what is called a filament.
\item \textbf{The all-pairs median is not a local-spacing estimator.} For points distributed uniformly along, or spaced regularly across, a segment of length $L$ it converges on $0.293\,L$; more generally it follows the extent of the occupied region rather than the local neighbour spacing, with a coefficient that depends on the distribution. The published HGBS values we could check are local or along-fibre statistics and are consistent with ours to within about a third in each cloud.
\item \textbf{The smallest separations are not interpretable.} The gap distributions truncate below $0.011$--$0.044$\,pc, entirely within the \textit{Herschel} beam and \textsc{getsources} deblending limit, so $S_{\rm local}$ is an upper limit on the true local separation of the gas fragments.
\item \textbf{The physical cause is undetermined.} A short intrinsic wavelength, fibre superposition, spatially varying stability and post-fragmentation evolution all remain viable and are not separable with column density alone.
\end{enumerate}

These results apply to four nearby, low-mass Gould Belt clouds of broadly solar metallicity and similar radiation field. Nothing here supports extrapolation to a general population, to high-mass complexes, or to different metallicity or radiation environments. The decisive next observations are velocity-resolved; the decisive theoretical question, whether an evolving filament must select the equilibrium wavelength at all, is taken up in Paper~II.

\section*{Acknowledgements}

We thank the \textit{Herschel} Gould Belt Survey team for the datasets.

\section*{Data availability}

The observational analysis pipeline is archived under a CC-BY-4.0 licence with the citable Zenodo DOI \texttt{10.5281/zenodo.14872301}, versioned to the accepted release. It contains the DisPerSE and FilFinder parameter sets, the bridging, association and adjacent-core-spacing scripts, the null-generation and bootstrap code, the derived per-component catalogues (source positions, arclengths and gaps, tabulated \emph{component by component} rather than pooled by cloud, so that every statistic reported here can be recomputed without re-running the extraction), and a dated commit history recording the order in which the analysis was carried out, including the point before any spacing statistic was computed at which the sample-selection criteria of Section~\ref{sec:distance_uncertainties} were fixed. The numerical material of Paper~II is deposited separately under \texttt{10.5281/zenodo.14872318}. HGBS column-density and skeleton maps are available from the \textit{Herschel} Science Archive, and the \textit{Gaia} DR3 distances from \citet{Zhang2023}.

\nocite{PaperII}
\bibliographystyle{mnras}

\appendix

\section{The limited HGBS sample}
\label{sec:limited_sample}
The four \emph{limited} regions fail at least one of the primary-selection criteria (Section~\ref{sec:distance_uncertainties}) and are excluded from the headline measurement. To show that the exclusion is not doing the work, we nevertheless ran two of them through the identical pipeline (Table~\ref{tab:limited}). Ophiuchus fails only on topology, and its heavily branched L1688 network gives the shortest values in the whole survey on both statistics. CRA fails only on core count, and falls inside the range spanned by the primary four. TMC1's deposited skeleton carries no usable astrometric header, and Serpens is excluded on distance and has no catalogue we use. Neither admissible cloud moves any statistic towards the equilibrium-cylinder reference. The deposited material lists all eight regions together with the core-count-weighted mean.

\begin{table}
\caption{The two \emph{limited} regions that can be passed through the unmodified pipeline, with the primary four repeated for comparison. Quality flag: \textsc{t}~= fails the topology criterion, \textsc{n}~= fails the core-count criterion. These values are reported for completeness and are not folded into any headline number, because the association step on a heavily branched network is not equivalent to the one performed on the primary sample.}
\label{tab:limited}
\footnotesize
\setlength{\tabcolsep}{4pt}
\begin{tabular}{lccccc}
\toprule
Region & flag & $N_{\rm cat}$ & $L_{\rm skel}$ (pc) & $S_{\rm local}$ & $S_{\rm skel}$ \\
\midrule
Ophiuchus & \textsc{t} & 513 & 32 & 0.6 & 0.8 \\
CRA       & \textsc{n} & 239 & 20 & 2.7 & 4.4 \\
\midrule
Orion B & --- & 1870 & 337 & 1.8 & 3.4 \\
Aquila  & --- &  749 & 106 & 1.7 & 3.4 \\
Perseus & --- &  816 & 290 & 3.0 & 5.8 \\
Taurus  & --- &  536 & 209 & 1.8 & 4.5 \\
\bottomrule
\end{tabular}
\end{table}




\section{The $L/3$ bias of the pairwise-median spacing}
\label{sec:pwmedian}
For $N$ points uniformly distributed on a filament of length $L$, a single pairwise distance $D=|X_i-X_j|$ has the triangular PDF $f(d)=2(L-d)/L^2$, with mean $\langle D\rangle=L/3$ and median $d^\star=L(1-1/\sqrt{2})\approx0.293\,L\approx0.88\,(L/3)$. The median of all $N(N-1)/2$ pairwise distances converges to $d^\star$ as $N\to\infty$, for a random or a perfectly periodic distribution alike, so the pairwise median measures the region extent, not the fragmentation wavelength. A Monte-Carlo over $N=5$--$500$ gives median$/L\approx0.293\,(1+0.43/N)$, i.e.\ $\le2\%$ inflation for $N\gtrsim20$, so the convergence is reached at the core counts of interest. For branched skeletons the statistic measures $\approx0.3\times$ the component's arclength extent, and for a realistic filament length distribution it tracks the rms-length$/3$ (full tables in the Zenodo archive). Computing the pairwise median over \emph{all} cores in each primary region (without per-filament segmentation) returns $\sim L_{\rm region}/3$ ($12.5$, $7.5$, $8.0$, $2.5$\,pc for Orion~B, Aquila, Perseus, Taurus), the region extent, confirming directly that the statistic measures segmentation scale, not the fragmentation wavelength.


\section{Systematics}
\label{sec:systematics}


Not all of the checks reported above carry equal weight, and it matters which are decisive. With $N=4$ clouds, the tests that bear directly on the conclusions are those that could change the \emph{sign} of a result: the segmentation comparison (Section~\ref{sec:ontology}), the coverage and association-radius sweeps, and the core-population splits (Section~\ref{sec:coverage}). The remainder, including the width and persistence conventions, the injection--recovery corrections and the eligible-length sweeps, are illustrative: they move numbers within the envelope already quoted and none of them reverses a conclusion. Table~\ref{tab:errorbudget} collects the terms under three headings, which are not commensurable and should not be combined. \emph{Measurement} terms would change under resampling of the same data. \emph{Calibration} terms are deterministic offsets from a fixed methodological choice. \emph{Object-definition} terms arise because different structural definitions identify different objects, and the skeleton algorithm belongs entirely to this third category: the factor of two between a persistent-crest and a medial-axis network is not an uncertainty bar on one physical quantity but a demonstration that the inferred point process depends on what is called a filament. A fourth heading collects \emph{modelling} terms internal to the simulations. Uncorrelated statistical terms are combined in quadrature, $\sigma_{\rm sys}^2=\sum_i\sigma_i^2$; the multiplicative width- and skeleton-convention factors are applied afterwards as separate analysis-choice envelopes and not added into that sum. The adopted $(s_{\rm ACS,med}/W)_0\approx1.4$ carries residual pipeline bias ($\sim10$ per cent), protostellar migration ($\sim7.5$ per cent), between-region scatter ($\sim4$ per cent) and a $\lesssim5$ per cent manual-editing proxy in quadrature, with the width convention ($\sim25$ per cent) and the skeleton algorithm (a factor $\sim2$) applied on top. Because the dominant term is a deterministic offset we quote a min--max envelope, $S_{\rm local}\approx0.8$--$1.9$ for DisPerSE, rather than a Gaussian $\sigma$.

\emph{Projection} acts against our conclusion and is shown in Table~\ref{tab:headline} for that reason. For an isotropic orientation distribution $s_{\rm obs}/s_{\rm true}=\sin i$ has median $0.87$, so deprojection raises the intrinsic spacing by $15$--$27$ per cent, moving the adopted $1.4$ to $1.6$--$1.8$. We do not apply it, because the orientation distribution within a single cloud need not be isotropic, but we carry it as a one-sided systematic.

\emph{Skeleton-algorithm dependence} is the largest observational term. The DisPerSE and FilFinder values ($\approx1.8$ and $\approx0.5$--$1.5$ in the fixed-width convention) bracket a genuine ambiguity in what constitutes ``the filament'', not a measurement error, and a matched-threshold FilFinder extraction shows the inferred spacing is itself strongly sensitive to the mask threshold, so the factor of two is if anything a lower bound. We base our quantitative conclusions on the persistent-crest network and treat the medial axis as a test of the structural definition rather than an equally weighted measurement. That preference is not purely conventional. On a synthetic map built by projecting a simulated filament, superposing a second at $30^\circ$, adding a smooth cirrus background and convolving with the beam, so that the true ridges are known exactly, the two constructions fail in opposite ways. The medial axis of a mask returns $2.3$--$4.5$ times the injected ridge length depending on the mask threshold, but only $12$--$30$ per cent of that length lies within one beam of a true ridge. The persistence-based crest extraction returns far less of the ridge, recovering under a tenth of it in this test, but $70$ per cent of what it does return lies within one beam, with a median offset of $0.1$ beams. The crest construction is thus incomplete and pure, the medial axis complete and impure, and the excess length of the latter is largely not on any real ridge. This is one synthetic realisation with a simplified stand-in for the published extraction, so it is indicative, not decisive; but it is the direction that matters here, because impurity inflates $L$ and therefore $S_{\rm skel}$, which is the sense in which the two observational networks differ.




\begin{table}
\caption{Terms on the measured $s_{\rm ACS,med}/W$ (symbols: Table~\ref{tab:principal}). Statistical and sensitivity terms are combined in quadrature; the width- and skeleton-convention factors are applied afterwards as a separate multiplicative \emph{analysis-choice envelope} and not as Gaussian uncertainties. The modelling terms that enter only when an observed spacing is matched to a simulated wavelength, and the numerical terms on the wavelength itself, belong to that comparison rather than to this measurement and are set out in Paper~II.}
\label{tab:errorbudget}
\footnotesize
\setlength{\tabcolsep}{3pt}
\begin{tabular}{@{}p{3.1cm}p{1.9cm}p{2.3cm}@{}}
\toprule
Source & Type & Magnitude \\
\midrule
Bootstrap / cloud-to-cloud & sensitivity & $\sigma\approx0.4$ ($\sim$20\%) \\
ACS pipeline bias & systematic & $\sim$10\% (corrected) \\
Protostellar migration & sensitivity & $\sim$7.5\% \\
Projection (deprojection) & systematic (one-sided, \emph{raises} $s$) & $+15$--$27\%$ \\
LOS blending (Aquila) & systematic (one-sided, \emph{shortens} $s$) & $\lesssim50\%$ worst case \\
Skeleton-coverage selection & systematic (one-sided, \emph{shortens} $s$) & $\lesssim17\%$ (bounded) \\
Width convention, fixed vs region & calibration & $\sim$25\% \\
Skeleton algorithm, crest vs medial axis & \emph{object definition} & factor $\sim$2 \\
Manual skeleton editing & \emph{unvalidated proxy bound} & $\lesssim$5\% \\
\bottomrule
\end{tabular}
\end{table}


\subsection*{Region-specific widths}

\begin{table}
\caption{Region-specific, beam-deconvolved filament widths from a direct Plummer fit to perpendicular column-density profiles along the skeletons (the HGBS width definition; \citealt{Arzoumanian2019}), and the resulting $S_{\rm local}=s_{\rm ACS,med}/W_{\rm fil}$. $W_{\rm fil}$ is the beam-deconvolved Plummer FWHM ($18.2''$ beam), fitted over a perpendicular half-range of $0.4$\,pc either side of the crest with the outer $0.1$\,pc used to set a local background, and $p$ is the Plummer index (a \emph{free} fit parameter, not fixed; it converges to $p\simeq2.00$ in all four regions, Taurus giving $2.01$, consistent with the canonical Plummer-2 profile), with $N$ the number of fitted profiles. Quoted $W_{\rm fil}$ uncertainties are the standard error on the median of the per-profile fits within each region and do not include the beam-deconvolution or background-subtraction systematics, which are carried separately in the error budget (Section~\ref{sec:systematics}); the fitted $p$ has a per-region dispersion of $\pm0.1$ across profiles. All four regions give $p\simeq2$ and place $S_{\rm local}$ below the reference band on the \emph{median} spacing quoted here; the \emph{mean} adjacent spacing is a factor $\approx1.3$--$2.2$ larger (Section~\ref{sec:pointprocess}), so the mean-based ratio is $\approx2.1$--$4.7$, still below the reference band in every region.}
\label{tab:widths}
\footnotesize
\setlength{\tabcolsep}{3pt}
\begin{tabular}{p{1.7cm}ccccc}
\toprule
Region & $W_{\rm fil}$ (pc) & $p$ & $N$ & $s_{\rm ACS,med}$ (pc) & $S_{\rm local}$ \\
\midrule
Taurus  & $0.110$ & $2.01$ & 295 & $0.181$ & $1.6$ \\
Orion B & $0.148$ & $2.00$ & 356 & $0.177$ & $1.2$ \\
Perseus & $0.147$ & $2.00$ & 235 & $0.297$ & $2.0$ \\
Aquila  & $0.161$ & $2.00$ & 446 & $0.166$ & $1.0$ \\
\midrule
median  & $0.148$ & $2.00$ & --- & $0.179$ & $1.4$ \\
\bottomrule
\end{tabular}
\par\vspace{2pt}
\footnotesize The measured widths ($0.11$--$0.16$\,pc) agree with the published HGBS values and exceed the canonical $0.10$\,pc, so region-specific Plummer widths give a \emph{lower} $S_{\rm local}$ (median $\approx1.4$, raw) than the raw fixed-width value ($\approx1.8$; Table~\ref{tab:headline}): all four regions are below the reference band under either convention. Adjacent-core spacings are the full-coverage medians (Section~\ref{sec:injection_recovery}), measured on the same DisPerSE skeleton as the widths (its deposited crest for the width fit, its bridged form for the spacing); the injection--recovery bias correction lowers each $S_{\rm local}$ by a further $\sim$10--20\%. (Pipeline outputs retained to 2 s.f.; the convention systematic dominates.)
\end{table}

\subsection*{Robustness to the persistence threshold}
\label{sec:persistence}

Since Section~\ref{sec:systematics} identifies the skeleton algorithm as the largest single systematic, the persistence threshold of our primary segmentation cannot be left untested. We cannot regenerate the published skeletons at a \emph{lower} threshold without the original DisPerSE run, so we approach the question from both sides.

\smallskip\noindent\emph{(i) Pruning the published skeleton upwards.} We decompose each skeleton into branches, terminating at junction pixels, and compute for each branch the persistence of its column-density maximum above the junction at which it attaches, in units of the local map noise $\sigma$ (estimated as $1.4826\times$MAD of the map minus a $5$-pixel median filter). Branches below a chosen threshold are removed and the procedure iterated to convergence. Raising the effective threshold from the published $3\sigma$ to $4\sigma$ removes $16$, $21$, $18$ and $15$ per cent of the skeleton length in Orion~B, Aquila, Perseus and Taurus, but changes the median adjacent spacing by only $-8.2$, $-2.7$, $+0.3$ and $-5.0$ per cent and the length per core by $-1.9$, $-3.2$, $-3.1$ and $-3.7$ per cent. Going further, to $6\sigma$ and $8\sigma$, changes nothing further at the per-cent level: the pruning saturates once the low-persistence spurs are gone. The occupancy fraction rises slightly, as expected when bare low-persistence branches are removed.

\smallskip\noindent\emph{(ii) An extraction with a tunable threshold.} Reaching the $2\sigma$ direction requires a crest extraction whose persistence threshold is a free parameter. We construct one directly from the column-density field: a descending-sweep merge tree pairs every maximum with the saddle at which it merges, maxima whose persistence falls below the chosen multiple of $\sigma$ are cancelled, and the surviving crests are traced by steepest-ascent walks from the retained saddles. This is not DisPerSE and does not reproduce the published skeleton pixel for pixel; it is used only to vary the threshold. Run at $2\sigma$, $3\sigma$ and $4\sigma$ on Orion~B it gives $S_{\rm local}=1.63$, $1.63$ and $1.63$ and $S_{\rm skel}=2.53$, $2.50$ and $2.48$, with the associated core count falling from $635$ to $597$ and the occupancy fraction rising from $0.41$ to $0.42$. The local statistic is thus insensitive to the persistence threshold at the third significant figure over a factor of two in threshold, and the global statistic moves by two per cent.

\smallskip\noindent Both tests agree: the persistence threshold is not what makes the DisPerSE and FilFinder answers differ. Between $2\sigma$ and $8\sigma$ the local spacing moves by less than ten per cent, whereas changing the structural definition from persistent crests to the medial axis of a mask moves it by a factor of two. The dominant systematic is therefore the \emph{definition} of a filament, not the tuning of a parameter within one definition, which is the conclusion of Section~\ref{sec:whatfilament}. The manual editing applied to the published skeletons cannot be varied in this way and remains an unquantified term.




\section{Sensitivity of the eligible-length diagnostic}
\label{sec:maskapp}

\begin{table}
\caption{Sensitivity of the core-masked eligibility diagnostic to the masking radius $r_{\rm mask}$. Cells give $S_{\rm elig}=(L_{\rm elig}/N_{\rm elig})/W$ with $N_{\rm elig}$ in brackets; $r_{\rm mask}=0$ is the uncorrected raw-map calculation and $0.06$\,pc is adopted. Orion~B ($1.56$--$2.0$) and Aquila ($3.14$--$3.21$) are comparatively stable; Perseus drifts and Taurus does not stabilise.}
\label{tab:maskrad}
\footnotesize
\setlength{\tabcolsep}{4pt}
\begin{tabular}{lcccc}
\toprule
$r_{\rm mask}$ (pc) & Orion~B & Aquila & Perseus & Taurus \\
\midrule
$0$ (raw) & 1.56 [417] & 3.14 [306] & 1.96 [199] & 0.81 [186] \\
$0.030$   & 1.89 [317] & 3.14 [306] & 3.68 [85]  & 1.80 [43] \\
$0.045$   & 2.04 [277] & 3.15 [305] & 3.91 [70]  & 1.61 [37] \\
$0.060$   & \textbf{2.02} [268] & \textbf{3.17} [303] & 4.13 [57] & --- [13] \\
$0.090$   & 1.98 [252] & 3.21 [296] & 4.45 [41]  & --- [1] \\
$0.120$   & 1.95 [238] & 3.20 [297] & 3.62 [37]  & --- [1] \\
\bottomrule
\end{tabular}
\end{table}

Table~\ref{tab:maskrad} gives the eligible line density of Section~\ref{sec:eligible} as a function of the radius within which column density is replaced around each associated core. Orion~B and Aquila are comparatively stable, varying by $\approx30$ and $\approx2$ per cent respectively over a factor of four in masking radius; Perseus drifts over $3.6$--$4.5$ while its eligible core count falls from $85$ to $37$; and Taurus does not stabilise at any radius. Only the first two clouds support the comparison.


\section{The column-density-conditioned null and its calibration}
\label{sec:nullspec}

\emph{The column-density-conditioned null.} The null of Sections~\ref{sec:occupancy} and \ref{sec:pointprocess} asks whether cores are clumped \emph{beyond} the clumping already present in the filament they lie on. It is an inhomogeneous Poisson process on the one-dimensional skeleton with intensity
\begin{equation}
\lambda(x)\;\propto\;\left[\max\!\left(N_{\rm H_2}(x)-N_0,\,0\right)\right]^{\alpha},
\label{eq:nullintensity}
\end{equation}
where $x$ is arclength, $N_0=10^{21}\,{\rm cm^{-2}}$ a floor preventing low-column pixels from dominating, and $\alpha$ a single exponent fitted per sample. The normalisation is fixed by $\int\lambda\,{\rm d}x=N_{\rm obs}$, so each realisation contains the observed number of cores. The exponent is fitted by maximum likelihood on the observed core positions; to avoid fitting and testing on the same data we also refit it under leave-one-cloud-out cross-validation, giving the $E_{\rm cv}$ column of Table~\ref{tab:percloud}. Fitted values are $\alpha=1.6$--$2.4$ and the conclusions are unchanged over $\alpha=1$--$3$. For each cloud we draw $N_{\rm MC}=10^4$ realisations. The reported probability is one-sided in the sense of excess clumping,
\begin{equation}
p=\frac{1+\#\{F_{\rm null}\geq F_{\rm obs}\}}{1+N_{\rm MC}},
\label{eq:pval}
\end{equation}
so a small $p$ would indicate a Fano factor larger than the null expects and a $p$ near unity the opposite. The pointwise $2.5$--$97.5$ percentile envelope is quoted for $g(s)$.

Two properties limit what the construction can do. It is conditional on the deposited skeleton, so it tests the arrangement of cores along a fixed network and not the network itself. And the column density entering Eq.~\ref{eq:nullintensity} is the observed map, which the cores themselves raise, so part of the signal under test is built into the model. Masking the cores and interpolating the background is the obvious remedy, and Section~\ref{sec:occupancy} reports the calibration that shows it does not work: populations drawn from the conditioned intensity, with the observed core imprint added, are flagged as over-dispersed at a nominal $p<0.05$ in $73$ per cent of trials in Taurus and $100$ per cent in Orion~B. The probabilities in Table~\ref{tab:percloud} are accordingly descriptive.


\section{Convention decoder}
\label{sec:decoder}

\begin{table}
\caption{Convention decoder for the \emph{conditional} local descriptor (symbols: Table~\ref{tab:notation}): read every observed $s_{\rm ACS,med}/W$ quoted in the text against it. This table tabulates the adjacent-core spacing, which is measured only within core-bearing stretches and is \emph{not} the quantity a one-fragment-per-wavelength theory predicts (that is $L/N$; Section~\ref{sec:twostats}). Values are given under each width/skeleton convention and correction state for the four primary HGBS regions (revised distances); every combination lies below the reference band (classical $\lambda_{\rm m}\approx5$--$6\times$FWHM). We adopt the directly-measured region-Plummer value as our reference, $\boldsymbol{(s_{\rm ACS,med}/W)_0\approx1.4}$; the fixed-width value ($\approx1.8$, retained for comparability with prior HGBS work) and all others are read against it. The ``periodic'' column applies the periodic injection--recovery correction to the raw value; the data-matched \emph{clustered} column (bracketed) applies $\div1.64$. The DisPerSE conventions fall below the reference band by a factor $\approx1.9$--$5.3$; across the full envelope the factor spans $\approx2$--$11$, the upper end set by the densest FilFinder skeleton.}
\label{tab:headline}
\footnotesize
\setlength{\tabcolsep}{4pt}
\begin{tabular}{p{4.6cm}cc}
\toprule
Quantity & raw & periodic \\
\midrule
\multicolumn{3}{@{}l}{\emph{Conditional local descriptor} $s_{\rm ACS,med}/W$ (core-bearing stretches only)}\\
Region Plummer, DisPerSE (adopted, $(s_{\rm ACS,med}/W)_0$) & $\mathbf{1.4}$ & $1.2$ \\
Fixed $0.10$\,pc, DisPerSE & $1.8$ & $1.6$ \\
FilFinder skeleton & $0.5$--$1.5$ & $0.5$--$1.5$ \\
Per-segment local width, $\Delta s/W_{\rm local}$ & \multicolumn{2}{c}{$0.6$--$1.1$} \\
\emph{Deprojected} (adopted row $\div\,\langle\sin i\rangle_{\rm med}=0.87$) & $1.6$ & $1.4$ \\
\emph{Deprojected} (fixed $0.10$\,pc) & $2.1$ & $1.8$ \\
\midrule
\multicolumn{3}{@{}l}{\emph{Unconditional skeleton line-density statistics} (the closest available proxy for what one-bead-per-wavelength theory predicts)}\\
Global length per core $(L/N)/W$, fixed $0.10$\,pc & \multicolumn{2}{c}{$3.4$--$5.8$} \\
Mean adjacent, fixed $0.10$\,pc (branch-jump inflated) & \multicolumn{2}{c}{$2.1$--$4.7$} \\
\midrule
\multicolumn{3}{@{}l}{\emph{Legacy comparator and reference scales}}\\
Pairwise median ($\approx0.293L$; biased, not a spacing) & $2.8$ & \\
Matched simulation (upper limit) & \multicolumn{2}{c}{$\lambda_{\rm turn}/W\lesssim1.0$--$1.3$} \\
Classical $\lambda_{\rm m}$ (F\&M 2012) & \multicolumn{2}{c}{$5$--$6$ ($=4\times$ diam.)} \\
\bottomrule
\end{tabular}
\\[2pt]
\footnotesize
\emph{Notes} The reference point is the region-Plummer ACS $(s_{\rm ACS,med}/W)_0\approx1.4$ (bold), one illustrative point within the convention range $0.5$--$1.8$; it is an adopted reference value, chosen for comparability with prior HGBS work, not a uniquely preferred physical scale (and not a formal confidence interval; Section~\ref{sec:systematics}). The data-matched \emph{clustered} branch ($\div1.64$ on the raw value, giving $\approx0.9$ region-Plummer, $\approx1.1$ fixed-width and $\approx0.3$--$0.9$ FilFinder) is \emph{not} shown as a column because it is not independent: it applies a correction derived by \emph{assuming} the observed clustering (Section~\ref{sec:injection_recovery}), so it is circular and serves only as a lower bound; we carry the raw value throughout. The fixed-width convention gives a \emph{higher} value ($\approx1.8$) because it uses $0.10$\,pc rather than the larger measured widths; FilFinder gives lower values still. The full envelope factor $\approx2$--$11$ (Section~\ref{sec:injection_recovery}) is $\lambda_{\rm m}/(s_{\rm ACS,med}/W)$: its lower endpoint $\approx2$ uses $(\lambda_{\rm m},s_{\rm ACS,med}/W)=(5.5W,3.0)$ (fixed-width DisPerSE, Perseus) and its upper endpoint $\approx11$ uses $(5.5W,0.49)$ (region-Plummer FilFinder, Taurus). We quote $s_{\rm ACS,med}/W$ to one decimal place; finer digits are not significant. The two \emph{deprojected} rows apply the median projection factor $\langle\sin i\rangle_{\rm med}=0.87$ for a \emph{random} distribution of filament orientations. That assumption need not hold within a single cloud complex, where a coherent magnetic field can align the filaments, so the correction is indicative, not exact (Section~\ref{sec:systematics}); this correction acts \emph{against} our conclusion, raising the ratio by $15$--$27$ per cent, and we show it in the table rather than only in the text so that the effect is visible in the same place as the adopted numbers. We nevertheless carry the uncorrected value as the headline, because deprojection assumes an isotropic orientation distribution that the filaments in a single cloud need not obey.
\end{table}

Table~\ref{tab:headline} decodes every $s_{\rm ACS}/W$ quoted in the main text against the width convention, skeleton algorithm and correction stage that produced it.

\begin{table}
\caption{Principal statistics for three nested core samples. ``All'' is the full published catalogue; ``starless$+$prestellar'' removes protostellar sources; ``prestellar'' retains only cores classified as prestellar. $N_{\rm assoc}$ is the number associated with a skeleton, CoV the coefficient of variation of the adjacent gaps. The Taurus prestellar row rests on $53$ associated cores and its $S_{\rm skel}$ should not be interpreted.}
\label{tab:popsplit}
\footnotesize
\setlength{\tabcolsep}{4pt}
\begin{tabular}{llcccc}
\toprule
Cloud & sample & $N_{\rm assoc}$ & $S_{\rm local}$ & $S_{\rm skel}$ & CoV \\
\midrule
Orion~B & all                  & 995 & 1.77 & 3.4 & 0.83 \\
        & starless$+$prestellar& 961 & 1.79 & 3.5 & 0.83 \\
        & prestellar           & 569 & 1.83 & 5.9 & 0.82 \\
\midrule
Aquila  & all                  & 308 & 1.66 & 3.4 & 0.65 \\
        & starless$+$prestellar& 267 & 1.78 & 4.0 & 0.62 \\
        & prestellar           & 223 & 1.78 & 4.8 & 0.73 \\
\midrule
Perseus & all                  & 501 & 2.97 & 5.8 & 0.96 \\
        & starless$+$prestellar& 438 & 2.99 & 6.6 & 0.95 \\
        & prestellar           & 282 & 2.83 & 10.3 & 1.01 \\
\midrule
Taurus  & all                  & 463 & 1.81 & 4.5 & 1.36 \\
        & starless$+$prestellar& 424 & 1.84 & 4.9 & 1.34 \\
        & prestellar           &  53 & 2.04 & (39.5) & 1.30 \\
\bottomrule
\end{tabular}
\end{table}

\begin{table*}
\caption{Comparison with published HGBS spacing measurements. Perseus and Taurus are the two regions in which our own gap sequence shows significant non-independence (runs test; Section~\ref{sec:pointprocess}) and in which the post-masking eligible-core samples collapse (Section~\ref{sec:eligible}); their entries should be compared with the published values with those caveats in mind.}
\label{tab:literature}
\setlength{\tabcolsep}{4pt}
\footnotesize
\setlength{\tabcolsep}{2.5pt}
\begin{tabular}{lcccccc}
\toprule
Region & ACS (pc) & Orig.\ HGBS & pairwise (pc) & Lit.\ (pc) & \textbf{Lit.\ statistic type} & Reference \\
\midrule
\textbf{Primary sample} &  &  &  &  &  &  \\
Orion B & \textbf{0.177} & 0.212 & 0.313 & 0.20--0.30 & local & \citet{Konyves2020} \\
Aquila  & \textbf{0.166} & 0.206 & 0.346 & 0.22--0.26 & local & \citet{Konyves2015} \\
Perseus & \textbf{0.297} & 0.199 & 0.248 & 0.22 & local & \citet{Pezzuto2021} \\
Taurus  & \textbf{0.181} & 0.206 & 0.198 & 0.19 & along-fibre & \citet{Hacar2013} \\
\midrule
\textbf{Limited sample} &  &  &  &  &  &  \\
Ophiuchus &, & 0.197 & 0.206 & 0.18 & local & \citet{Konyves2020} \\
Serpens  &, & 0.188 & 0.331 & 0.17--0.19 & local & \citet{Konyves2020} \\
TMC1     &, & 0.203 & 0.195 & $\sim$0.20 & along-fibre & \citet{Hacar2013} \\
CRA      &, & 0.212 & 0.248 & $\sim$0.21 & local & \citet{Konyves2020} \\
\bottomrule
\end{tabular}

\vspace{0.1in}
\footnotesize
\emph{Notes} $^{a}$ The ``Lit.\ statistic'' column records what kind of statistic each cited value represents. All are \emph{local} separation measures, either nearest-neighbour or along-fibre adjacent separations, rather than all-pairs medians, and that is precisely why they agree with our ACS column and not with our pairwise column; the exact definitions differ in detail between papers. The region-extent bias identified here therefore applies wherever an all-pairs or pooled separation statistic is adopted, including the comparability statistic carried in this paper and in cross-region compilations, and not to these particular published values. $^{b}$ The primary statistic is the \emph{adjacent-core} (ACS) spacing (bold); the ``pairwise'' column is a legacy comparator only, inflated towards $L/3$ (Section~\ref{sec:statistics}) and not a physical spacing. ACS spacings for the limited sample are omitted. Perseus is the one region where the ACS ($0.297$\,pc) exceeds the pairwise median ($0.248$\,pc), because arclength ordering on its heavily bridged skeleton inflates path distances; we verified this is not a bridging artefact, an independent minimum-spanning-tree ordering of the same cores gives a consistent value, insensitive to the morphological bridging, and both remain below the reference band. $^{c}$ ``Original HGBS'' values use pre-Gaia distances; literature ranges reflect different methods and distance assumptions. $^{d}$ For Orion~B the pre-Gaia distance was already $\sim$400\,pc \citep{Konyves2020}, so the pairwise $0.212\to0.313$\,pc difference is a statistic-definition effect, not a distance revision.
\end{table*}

\section{Comparison with the literature}
\label{sec:lit}

Our revised-distance ACS measurements (Table~\ref{tab:literature}) differ from the original HGBS values by the region-dependent distance revisions; the core-count-weighted primary-region mean rises only $\sim10$ per cent, a change dominated by Aquila. Toggling each methodological choice on the Orion~B data at fixed distance shows that the \emph{statistic definition} dominates: switching from ACS to the pairwise median inflates the spacing from $0.177$ to $0.313$\,pc, a factor of $1.8$, and accounts for essentially the whole apparent shift, while skeleton re-extraction, association half-width, arclength ordering and the distance revision are each $\leq15$ per cent and largely cancel.

The published values we could check are local or along-fibre separations, not all-pairs medians, and our adjacent-core spacings scatter about them without a systematic offset: Taurus lies $5$ per cent below the along-fibre value of \citet{Hacar2013}, Orion~B and Aquila $12$ and $25$ per cent below the lower ends of the quoted local ranges, and Perseus $35$ per cent above. That is why our measurements are \emph{not} in tension with earlier HGBS work despite reporting much shorter values than the historical ``characteristic spacing'' narrative implies: the consistency is with local and along-fibre statistics, and the disagreement is with the pooled all-pairs statistic, which exceeds our adjacent-core value by factors of $1.1$--$2.1$ in the same four clouds. The region-extent bias applies wherever an all-pairs or pooled statistic is adopted, including the comparability statistic carried here and in cross-region compilations, and we make no claim that it affects these particular published values.




\bibliography{references_complete}

\bsp
\label{lastpage}
\end{document}
